\documentclass[11pt, a4paper, UTF8]{ctexart}

% --- 宏包加载 ---
\usepackage{amsmath, amssymb, amsthm} % 数学公式与符号
\usepackage{geometry}                % 页边距设置
\usepackage{enumitem}                % 列表格式自定义
\usepackage{titlesec}                % 标题格式自定义
\usepackage{hyperref}                % 超链接与目录
\usepackage{xcolor}                  % 颜色设置
\usepackage{bm}                      % 粗体矢量
%%这是验证修改的注释
% --- 页面设置 ---
\geometry{left=2.5cm, right=2.5cm, top=3cm, bottom=3cm}
\hypersetup{
    colorlinks=true,
    linkcolor=blue,
    urlcolor=cyan,
    pdftitle={有限元理论内功修炼路线图}
}

% --- 自定义命令 ---
\newcommand{\weakform}[1]{\int_{\Omega} #1 \, d\Omega}
\newcommand{\grad}{\nabla}

% --- 标题信息 ---
\title{\textbf{有限元理论内功修炼路线图：从代码实现回归数学本源}}
%\author{力学博士研究生理论补强计划}
\date{\today}

\begin{document}

\maketitle
%%
\begin{abstract}
    本路线图专为具备有限元编程经验（如六面体单元、动力学求解）但缺乏严谨理论推导基础的力学专业博士生设计。目标是打通从“强形式”微分方程到“弱形式”积分方程，再到离散矩阵方程的逻辑链路，最终实现理论与代码的深度耦合。
\end{abstract}

\tableofcontents

\newpage

\section{第一章：数学语言与连续介质基础}
在进入有限元之前，必须统一符号体系，建立严谨的数学基础。本章的目标是让你具备无障碍推导连续介质力学与有限元控制方程的能力。以下内容构成后续所有弱形式、虚功原理与离散推导的语言基础。

\subsection{指标记号与张量分析}
力学推导中，向量和张量的矩阵写法往往在求导时显得笨拙，指标符号是必须要掌握的“外语”。

\textbf{1. 爱因斯坦求和约定 (Einstein Summation Convention):}
\begin{itemize}
    \item \textbf{规则：} 在同一项中，若某个下标重复出现两次（称为哑指标 dummy index），则表示对该指标在 $1,2,3$（三维空间）上求和。只出现一次的下标称为自由指标 (free index)。
    \item \textbf{示例：} 
    向量点乘：$\bm{u} \cdot \bm{v} = u_i v_i = u_1 v_1 + u_2 v_2 + u_3 v_3$。
    矩阵乘向量：$\bm{y} = \bm{A} \bm{x}$ 写为 $y_i = A_{ij} x_j$。
    矩阵乘法：$\bm{C} = \bm{A} \bm{B}$ 写为 $C_{ij} = A_{ik} B_{kj}$。
\end{itemize}

\textbf{2. Kronecker 符号 $\delta_{ij}$ (克罗内克 delta):}
\begin{itemize}
    \item \textbf{定义：} 当 $i=j$ 时 $\delta_{ij}=1$；当 $i \neq j$ 时 $\delta_{ij}=0$。它相当于单位矩阵 $\bm{I}$ 的元素。
    \item \textbf{性质（替换作用）：} $\delta_{ij} u_j = u_i$（把 $j$ 强行换成了 $i$）。$\delta_{ii} = \delta_{11} + \delta_{22} + \delta_{33} = 3$（三维空间下单位矩阵的迹）。
\end{itemize}

\textbf{3. 置换符号 $e_{ijk}$ (Levi-Civita Symbol):}
\begin{itemize}
    \item \textbf{定义：} 当 $i,j,k$ 构成正置换 (123, 231, 312) 时值为 1；负置换 (132, 213, 321) 时值为 -1；有任意两指标相同时值为 0。
    \item \textbf{应用（叉乘）：} $\bm{c} = \bm{a} \times \bm{b}$ 可以写为 $c_i = e_{ijk} a_j b_k$。
    \item \textbf{重要的恒等式（$e-\delta$ 恒等式）：} $e_{ijk} e_{ilm} = \delta_{jl}\delta_{km} - \delta_{jm}\delta_{kl}$。这在推导诸如 $\nabla \times (\nabla \times \bm{u})$ 的公式时是核心工具。
\end{itemize}

\textbf{4. 梯度算子 $\nabla$ 与指标形式：}
\begin{itemize}
    \item $\nabla = \left( \frac{\partial}{\partial x_1}, \frac{\partial}{\partial x_2}, \frac{\partial}{\partial x_3} \right)$，用指标表示为 $\partial_i$ 或 $\frac{\partial}{\partial x_i}$。有时缩写为逗号下标：$u_{i,j} = \frac{\partial u_i}{\partial x_j}$。
    \item 散度 $\nabla \cdot \bm{v} = \frac{\partial v_i}{\partial x_i} = v_{i,i}$。
    \item 旋度 $\nabla \times \bm{v}$ 的第 $i$ 个分量是 $e_{ijk} \partial_j v_k$。
\end{itemize}

\subsection{高斯散度定理}
散度定理是连接积分域内部与边界的桥梁，也是整个有限元分部积分与推导弱形式的绝对核心。理解它在不同维度下的表现形式至关重要。

\textbf{1. 三维空间下的标准形式 (3D Divergence Theorem):}
\begin{itemize}
    \item \textbf{物理图像：} 区域 $\Omega$ 内所有“源”的总和（散度的体积分），等于通过包围该区域的闭合曲面 $\partial\Omega$ 流向外部的总流量（法向通量的面积分）。
    \item \textbf{向量形式：} 
    \[ \iiint_{\Omega} (\nabla \cdot \bm{v}) \, dV = \iint_{\partial\Omega} (\bm{v} \cdot \bm{n}) \, dS \]
    其中 $\bm{n} = (n_x, n_y, n_z)$ 是曲面外法向单位向量，$dS$ 是面积微元。
    \item \textbf{指标形式：} $\int_{\Omega} v_{i,i} \, d\Omega = \int_{\partial\Omega} v_i n_i \, d\Gamma$ （$i=1,2,3$）
\end{itemize}

\textbf{【三维例题：验证单位立方体上的散度定理】}

设向量场
\[
\bm{v}(x,y,z) = (x,\, y,\, z),
\]
积分区域取单位立方体
\[
\Omega = [0,1]\times[0,1]\times[0,1].
\]
我们分别计算体积分 $\iiint_{\Omega} \nabla\cdot\bm{v}\, dV$ 与边界通量积分 $\iint_{\partial\Omega} \bm{v}\cdot\bm{n}\, dS$，看看它们是否相等。

\textbf{第一步：先算散度。}
\[
\nabla \cdot \bm{v}
=
\frac{\partial x}{\partial x}
+
\frac{\partial y}{\partial y}
+
\frac{\partial z}{\partial z}
= 1+1+1=3.
\]
所以体积分变成
\[
\iiint_{\Omega} (\nabla \cdot \bm{v})\, dV
=
\iiint_{\Omega} 3\, dV
=
3 \cdot \text{Vol}(\Omega)
=
3 \times 1 = 3.
\]

\textbf{第二步：把闭合曲面拆成六个面逐个计算通量。}
单位立方体的边界由六个平面组成：
\[
x=0,\quad x=1,\quad y=0,\quad y=1,\quad z=0,\quad z=1.
\]
每个面都要取\textbf{外法向}。

\textbf{(1) 面 $x=1$：}
外法向为 $\bm{n}=(1,0,0)$，故
\[
\bm{v}\cdot\bm{n} = (x,y,z)\cdot(1,0,0)=x=1.
\]
该面的面积微元为 $dS=dy\,dz$，积分区域是 $(y,z)\in[0,1]^2$，所以
\[
\iint_{x=1} \bm{v}\cdot\bm{n}\, dS
=
\int_0^1\int_0^1 1\, dy\, dz
=1.
\]

\textbf{(2) 面 $x=0$：}
外法向为 $\bm{n}=(-1,0,0)$，故
\[
\bm{v}\cdot\bm{n} = (x,y,z)\cdot(-1,0,0)=-x=0.
\]
所以该面的通量为 0。

\textbf{(3) 面 $y=1$：}
外法向为 $\bm{n}=(0,1,0)$，故
\[
\bm{v}\cdot\bm{n}=y=1,
\]
因此
\[
\iint_{y=1} \bm{v}\cdot\bm{n}\, dS = 1.
\]

\textbf{(4) 面 $y=0$：}
外法向为 $\bm{n}=(0,-1,0)$，故
\[
\bm{v}\cdot\bm{n}=-y=0,
\]
因此该面的通量为 0。

\textbf{(5) 面 $z=1$：}
外法向为 $\bm{n}=(0,0,1)$，故
\[
\bm{v}\cdot\bm{n}=z=1,
\]
因此
\[
\iint_{z=1} \bm{v}\cdot\bm{n}\, dS = 1.
\]

\textbf{(6) 面 $z=0$：}
外法向为 $\bm{n}=(0,0,-1)$，故
\[
\bm{v}\cdot\bm{n}=-z=0,
\]
因此该面的通量为 0。

\textbf{第三步：把六个面的结果加总。}
\[
\iint_{\partial\Omega} \bm{v}\cdot\bm{n}\, dS
=
1+0+1+0+1+0 = 3.
\]

\textbf{结论：}
\[
\iiint_{\Omega} (\nabla\cdot\bm{v})\, dV
=
\iint_{\partial\Omega} \bm{v}\cdot\bm{n}\, dS
=3.
\]
这就把三维散度定理完整验证了一遍。

\textit{这个例题最值得体会的地方是：} 体内均匀存在散度 $3$，意味着区域内处处都像有“源”在向外吐流量；而边界通量积分告诉你，这些内部产生的总流量，最终完全通过外表面流了出去。有限元中的分部积分，本质上正是在做这种“体内导数 $\rightarrow$ 边界通量”的转换。

\textbf{2. 降维打击：二维空间与格林公式 (2D \& Green's Theorem):}
\begin{itemize}
    \item 当问题退化为平面问题（平面应力/应变）时，散度定理退化为著名的二维格林公式。
    \item \textbf{公式表述：} 区域内的双重积分等于其闭合边界曲线上的线积分：
    \[ \iint_{A} \left( \frac{\partial v_x}{\partial x} + \frac{\partial v_y}{\partial y} \right) dx dy = \oint_{C} (v_x n_x + v_y n_y) ds \]
    这里的 $\bm{n} = (n_x, n_y)$ 是曲线 $C$ 的单位外法向量，$ds$ 是弧长微元。
    \item \textbf{与传统格林公式的联系：} 在高数中你可能学过 $\iint_A (\frac{\partial Q}{\partial x} - \frac{\partial P}{\partial y}) dxdy = \oint_C (Pdx + Qdy)$。其实它们完全等价！因为曲线的单位法向量 $\bm{n} ds = (dy, -dx)$。如果你令 $v_x = Q, v_y = -P$，上述散度定理立刻变成传统格林公式。在力学和有限元推导中，我们更喜欢用法向量 $\bm{n}$ 的形式，因为它具有更明确的力学意义（如面力方向）。
\end{itemize}

\textbf{【二维例题：验证单位正方形上的格林公式】}

设二维向量场
\[
\bm{v}(x,y) = (x^2,\, y),
\]
区域取单位正方形
\[
A = [0,1]\times[0,1].
\]
我们验证
\[
\iint_A \left(\frac{\partial v_x}{\partial x}+\frac{\partial v_y}{\partial y}\right)\, dxdy
=
\oint_C (v_x n_x + v_y n_y)\, ds.
\]

\textbf{第一步：计算区域内部的散度积分。}
因为
\[
\frac{\partial v_x}{\partial x} = \frac{\partial x^2}{\partial x}=2x,
\qquad
\frac{\partial v_y}{\partial y} = \frac{\partial y}{\partial y}=1,
\]
所以
\[
\nabla\cdot\bm{v}=2x+1.
\]
于是
\[
\iint_A (2x+1)\, dxdy
=
\int_0^1\int_0^1 (2x+1)\, dx\, dy.
\]
先对 $x$ 积分：
\[
\int_0^1 (2x+1)\, dx = \left[x^2+x\right]_0^1 = 2.
\]
再对 $y$ 积分：
\[
\int_0^1 2\, dy = 2.
\]
因此体内积分结果为
\[
\iint_A \nabla\cdot\bm{v}\, dA = 2.
\]

\textbf{第二步：把边界 $C$ 分成四条边分别计算线积分。}
单位正方形的边界可拆成
\[
x=0,\quad x=1,\quad y=0,\quad y=1.
\]
注意这里不是面积分，而是沿边界的\textbf{线积分}，但 integrand 仍然是 $\bm{v}\cdot\bm{n}$。

\textbf{(1) 右边 $x=1$：}
外法向为 $\bm{n}=(1,0)$，弧长微元为 $ds=dy$，故
\[
\bm{v}\cdot\bm{n} = (x^2,y)\cdot(1,0)=x^2=1.
\]
因此
\[
\int_{x=1} \bm{v}\cdot\bm{n}\, ds
=
\int_0^1 1\, dy = 1.
\]

\textbf{(2) 左边 $x=0$：}
外法向为 $\bm{n}=(-1,0)$，故
\[
\bm{v}\cdot\bm{n} = (x^2,y)\cdot(-1,0) = -x^2 = 0,
\]
因此该边贡献为 0。

\textbf{(3) 上边 $y=1$：}
外法向为 $\bm{n}=(0,1)$，弧长微元为 $ds=dx$，故
\[
\bm{v}\cdot\bm{n} = (x^2,y)\cdot(0,1)=y=1.
\]
因此
\[
\int_{y=1} \bm{v}\cdot\bm{n}\, ds
=
\int_0^1 1\, dx = 1.
\]

\textbf{(4) 下边 $y=0$：}
外法向为 $\bm{n}=(0,-1)$，故
\[
\bm{v}\cdot\bm{n} = (x^2,y)\cdot(0,-1)=-y=0,
\]
因此该边贡献为 0。

\textbf{第三步：汇总边界通量。}
\[
\oint_C \bm{v}\cdot\bm{n}\, ds
=
1+0+1+0 = 2.
\]

\textbf{结论：}
\[
\iint_A \left(\frac{\partial v_x}{\partial x}+\frac{\partial v_y}{\partial y}\right)\, dxdy
=
\oint_C (v_x n_x + v_y n_y)\, ds
=2.
\]
这就验证了二维格林公式。

\textit{这个二维例题的关键启发是：} 在二维有限元里，域内的散度项经过分部积分后，会转化成边界上的法向通量项。也就是说，边界边上的力、热流或通量，绝不是“凭空加上去”的，而是从这个二维版本的散度定理中自然长出来的。

\textbf{3. 一维空间与牛顿-莱布尼茨公式 (1D \& Fundamental Theorem of Calculus):}
\begin{itemize}
    \item 在一维线单元中，区域积分就是定积分，边界就是区间的两个端点。
    \item \textbf{公式表述：} 
    \[ \int_{x_1}^{x_2} \frac{dv}{dx} \, dx = v(x_2) - v(x_1) \]
    这可以看作一维的“散度定理”，其中右端点 $x_2$ 的外法向是向右（$n=+1$），左端点 $x_1$ 的外法向是向左（$n=-1$），所以通过边界的通量就是 $v(x_2)(+1) + v(x_1)(-1)$。
\end{itemize}

\textbf{4. 其他重要推广形式（记住核心：求导 $\partial_i$ $\rightarrow$ 法向 $n_i$）：}
只要掌握了“把偏导数 $\partial_i$ 转化为边界法向量 $n_i$”这个核心思想，散度定理有许多变体，在连续介质力学中极为常用：
\begin{itemize}
    \item \textbf{标量梯度的积分：} $\int_{\Omega} \nabla \phi \, d\Omega = \int_{\partial\Omega} \phi \bm{n} \, d\Gamma$。
    （指标证明思路：将算子换位，$\int_{\Omega} \phi_{,i} \, d\Omega = \int_{\partial\Omega} \phi n_i \, d\Gamma$）。
    \item \textbf{向量旋度的积分：} $\int_{\Omega} \nabla \times \bm{v} \, d\Omega = \int_{\partial\Omega} \bm{n} \times \bm{v} \, d\Gamma$。
    \item \textbf{二阶张量（如应力）的积分：} 这是力学中最常用的形式，将体内平衡转换到边界上面力：
    \[ \int_{\Omega} \nabla \cdot \bm{\sigma} \, d\Omega = \int_{\partial\Omega} \bm{\sigma} \cdot \bm{n} \, d\Gamma \quad \text{或指标形式：} \quad \int_{\Omega} \sigma_{ij,j} \, d\Omega = \int_{\partial\Omega} \sigma_{ij} n_j \, d\Gamma \]
    右侧 $\sigma_{ij} n_j$ 正是柯西应力公式在边界上定义的表面力向量 $t_i$。
\end{itemize}

\subsection{分部积分公式}
如果你刚上大一，刚刚学过微积分，你一定背过这个公式：$\int u dv = uv - \int v du$。
在有限元理论中，分部积分法是绝对的“C位魔法”，它的核心作用只有两个字：\textbf{“甩锅”}。

\textbf{1. 从大一的高数说起（一维的情况）：}
\begin{itemize}
    \item 假设你要算积分 $\int_{a}^{b} u(x) \cdot v''(x) \, dx$（注意 $v$ 上面有两撇，也就是 $v$ 的二阶导数）。
    \item 这个积分很让人头疼，因为它要求 $v(x)$ 必须光滑到能够求两次导数（也就是二阶连续可导），条件非常苛刻！
    \item 怎么办？我们用高数里的分部积分公式“甩锅”：
    \[ \int_{a}^{b} u \cdot v'' \, dx = \left[ u \cdot v' \right]_a^b - \int_{a}^{b} u' \cdot v' \, dx \]
    \item \textbf{魔法发生了：} 右边的积分项里，变成了 $u'$ 和 $v'$ 的乘积。我们成功地把 $v$ 脑袋上的一个“导数撇号（$'$）”，**甩给**了 $u$！
    \item \textbf{结果：} 现在，不需要 $v$ 有二阶导数了，只要 $u$ 和 $v$ 都有\textbf{一阶导数}就能算这个积分。这就是有限元里常说的“\textbf{弱化了对试函数的连续性要求}”（从强形式变成了弱形式）。同时，公式产生的那个边界项 $\left[ u \cdot v' \right]_a^b$，正好用来代入我们在边界上受到的力（自然边界条件）！
\end{itemize}

\textbf{2. 升维！三维空间的分部积分：}
在三维空间中，求导从一个撇号 $'$ 变成了梯度算子 $\nabla$。我们怎么把 $\nabla$ “甩锅”呢？这要借助前面刚学的\textbf{散度定理}！
\begin{itemize}
    \item \textbf{第一步：乘积求导法则（跟一维 $(uv)'=u'v+uv'$ 完全一样）。} 
    考察两个标量函数 $u$ 和 $v$。算子 $\nabla$ 作用在乘积 $u \nabla v$ 上：
    \[ \nabla \cdot (u \nabla v) = (\nabla u) \cdot (\nabla v) + u (\nabla \cdot \nabla v) \]
    整理一下，把带有“二阶导数”（拉普拉斯算子 $\nabla \cdot \nabla v = \nabla^2 v$）的项留在左边：
    \[ u \nabla^2 v = \nabla \cdot (u \nabla v) - (\nabla u) \cdot (\nabla v) \]
    \item \textbf{第二步：两边在体积 $\Omega$ 上求积分。}
    \[ \int_{\Omega} u \nabla^2 v \, d\Omega = \int_{\Omega} \nabla \cdot (u \nabla v) \, d\Omega - \int_{\Omega} (\nabla u) \cdot (\nabla v) \, d\Omega \]
    \item \textbf{第三步：散度定理出场！}
    看等号右边的第一项，它是对一个散度 $\nabla \cdot (\dots)$ 求体积分！根据散度定理（体积分转化为面积分），它直接变成了在边界 $\partial\Omega$ 上的积分：
    \[ \int_{\Omega} \nabla \cdot (u \nabla v) \, d\Omega = \int_{\partial\Omega} (u \nabla v) \cdot \bm{n} \, d\Gamma \]
    \item \textbf{第四步：得到最终的“甩锅”公式。} 将上一步代入，得到三维分部积分标准公式：
    \[ \int_{\Omega} u \nabla^2 v \, d\Omega = \int_{\partial\Omega} u (\nabla v \cdot \bm{n}) \, d\Gamma - \int_{\Omega} (\nabla u) \cdot (\nabla v) \, d\Omega \]
    你看，左边 $v$ 头顶上的二阶导数 $\nabla^2$，在右边成功变成了 $u$ 和 $v$ 各自只承担一个一阶导数 $\nabla$！边界项中的 $\nabla v \cdot \bm{n}$ 通常代表边界上的法向导数（通量）。
\end{itemize}

\textbf{3. 用指标符号再看一次（为了装得像个内行）：}
如果你看懂了上面，我们用张量指标符号写一遍，极其简洁清爽：
\begin{itemize}
    \item 乘积求导：$(u v_{i,i}) = (u v_i)_{,i} - u_{,i} v_i$
    \item 两边积分并对中间项用散度定理：
    \[ \int_{\Omega} u v_{i,i} \, d\Omega = \int_{\partial\Omega} u v_i n_i \, d\Gamma - \int_{\Omega} u_{,i} v_i \, d\Omega \]
\end{itemize}

\textbf{总结：} \textbf{分部积分就是能量转换的魔术。} 它把方程内部严苛的求导要求“降阶”分摊给权函数，并“吐”出了一个边界积分项。在第七章的“泊松方程推导”中，你会亲自上手体验这个过程。

\textbf{第一章练习建议：}
尝试纯手工推导证明张量恒等式 $\nabla \times (\nabla \phi) = \bm{0}$ 和 $\nabla \cdot (\nabla \times \bm{u}) = 0$。体会指标符号 $e_{ijk} \partial_j \partial_k \phi = 0$ 中因为 $e_{ijk}$ 反对称和偏导数对称导致的抵消。若能熟练推导，说明指标符号已过关。

\section{第二章：加权余量法与 Galerkin 思想}
如果说变分法是物理学家的“内功”（通过物理能量极值推导），那么加权余量法就是纯粹的数学“外功”。它能让你在不知道系统的“能量”到底是什么的情况下，强行把任何偏微分方程变成有限元能求解的积分形式。

\subsection{强形式与残差}
\begin{itemize}
    \item \textbf{强形式 (Strong Form)：}
    大自然给我们的物理定律（比如流体力学的 N-S 方程，固体力学的平衡方程），往往写成偏微分方程的形式：
    \[ \mathcal{L}(u) - f = 0 \quad (\text{比如 } \mathcal{L} \text{ 是 } -\nabla^2, f \text{ 是热源}) \]
    为什么叫它“强”？因为它太“霸道”了！它要求这个等式在区域 $\Omega$ 内的\textbf{每一个几何微观点}上都必须绝对等于 0，且要求解 $u$ 高阶可导。
    
    \item \textbf{残差 (Residual)：}
    有限元实际上是用多项式（形函数 $N_i$）去拼凑、去猜真实解 $u$。我们把猜出来的近似解记作 $\tilde{u}$。
    既然是拼凑猜出来的，把它强行代入“强形式”方程里，方程左边显然不可能完美等于 0。这个“差一点等于 0”的误差，就叫\textbf{残差 $R$}：
    \[ R = \mathcal{L}(\tilde{u}) - f \neq 0 \]
    我们的目标，就是想办法让这个残差 $R$ 尽可能地逼近于 0。
\end{itemize}

\subsection{加权积分与局部支集}
既然我们找不到一个能在所有点上都让残差为 0 的神仙解 $\tilde{u}$，我们就只能开始“妥协”。
\begin{itemize}
    \item \textbf{全域平均：} 我们不再要求 $R$ 在每个点都为 0，只要它在整个区域 $\Omega$ 内的“总体平均效应”为 0 就行了。
    但是，如果只是简单地求个积分 $\int R \, d\Omega = 0$，可能会出现“东边正的误差很大，西边负的误差很大，加起来刚好抵消”这种自欺欺人的情况。
    \item \textbf{引入权函数 (Weighting Function) 实行“分片包干”：}
    为了避免“东边的巨坑填了西边的高山”这种正负胡乱抵消，数学家找来了一大批“质检员”（称为权函数 $w_1, w_2, \dots, w_n$）。
    这些质检员有一个绝活：\textbf{每个人只负责盯住一块特定的极小区域（这叫局部支集 Local Support）}，在其他地方他们的值是 0。
    
    \textbf{【举个通俗的例子让你秒懂】：}
    假设我们有一根 10 米长的梁，我们把残差 $R(x)$ 想象成梁上各处的“施工误差”。
    如果不用权函数，直接算总误差 $\int_0^{10} R(x) dx$。如果第 1 米处误差是 $+100$（多建了），第 9 米处误差是 $-100$（少建了），加起来总误差是 0，工程队说“完美过关”，但这显然是荒谬的！
    
    现在，我们派 11 个质检员 $w_0, w_1, \dots, w_{10}$ 站在 $x=0, 1, \dots, 10$ 米的地方。
    \begin{itemize}
        \item 质检员 $w_1$（站在 1 米处），他手里的手电筒只能照亮他脚下的 $[0, 2]$ 米区域（这就是他的局部支集）。在 2 米之外，他的光照不到（权值为 0）。
        \item 我们强硬地要求：误差在 $w_1$ 的光照下，积分必须为 0，即 $\int_0^{10} w_1(x) \cdot R(x) dx = \int_0^2 w_1(x) \cdot R(x) dx = 0$。
        \item 此时，第 1 米处的巨大正误差 $+100$ 暴露在 $w_1$ 的眼皮底下，而第 9 米处的负误差 $-100$ 根本不在他的光照范围内！
        \item 所以，$w_1$ 这个方程绝不可能靠第 9 米的误差来抵消，它只能逼迫**第 1 米附近的误差老老实实地趋近于 0**。
    \end{itemize}
    当所有的质检员（成百上千个权函数方程）手电筒的光圈密密麻麻地覆盖并互相重叠时，整个区域的误差就被死死地“锁”在了极小的片区内，没有任何胡乱抵消的余地。这就是大名鼎鼎的\textbf{加权余量法}！
    
    \vspace{0.5em}
    \textbf{【灵魂拷问：具体怎么实现“质检员”只监管一小片区域？】}
    
    你可能会问：在数学和代码上，到底是用什么机制，硬生生地把一个函数“裁剪”到只在一小块区域有值的？答案就在于\textbf{网格划分 (Meshing)} 与 \textbf{分段多项式 (Piecewise Polynomials)}。
    \begin{itemize}
        \item \textbf{数学上的“帽子函数”：} 我们人为构造了一种极其特殊的分段函数（有限元中称为形函数 $N_i$）。以一维为例，质检员 $w_i$（对应节点 $i$）的权函数就像一个搭在节点上的“尖顶小帐篷”：在自己脚下节点处值为 1，向两侧相邻节点线性衰减到 0，\textbf{而在帐篷以外的所有区域，强行规定其值严格为 0！} 出了帐篷，任何误差乘上 0 还是 0，质检员自然就“看不见”了。这在数学上叫作“紧支集 (Compact Support)”。
        \item \textbf{代码实现的反向思维：} 在写代码时，我们不可能写 1000 个复杂的 \texttt{if-else} 来判断当前坐标是否在某个帐篷里。有限元的程序设计采用了极其巧妙的\textbf{以单元（Element）为中心}的视角。代码的本质是一个外层循环（\texttt{for elem in elements:}）。当程序进入某一个具体的单元时，它只在内部激活与该单元相连的少数几个局部节点的权函数；对于全局中其他所有遥远的权函数，程序逻辑天然认为它们在这个单元内的值为 0，根本不让它们参与当前单元的积分计算。算完这个局部小矩阵后，再根据节点编号将其像贴瓷砖一样“组装 (Assembly)”到全局大矩阵中。
    \end{itemize}

    \item \textit{旁支小知识：} 你选择什么样的“质检员”（权函数 $w$），就会得到不同的数值方法。比如，如果质检员只在某个点盯着看（用狄拉克 $\delta$ 函数），这就是配点法 (Collocation)；如果质检员只管一块小区域（阶跃函数），这就是子域法或有限体积法 (FVM)。
\end{itemize}

\subsection{Galerkin 方法}
在众多的“质检员”候选人中，有一个选择脱颖而出，成为了现代有限元（尤其是固体力学）的绝对核心，这就是 \textbf{Galerkin 法 (伽辽金法)}。

\begin{itemize}
    \item \textbf{核心思想：让“质检员”和“施工队”同源！}
    我们的近似解 $\tilde{u}$ 是由基函数（形函数）$N_j$ 作为“施工队”拼出来的：$\tilde{u} = \sum N_j u_j$。
    Galerkin 的绝妙主意是直接规定：\textbf{我们就拿施工用的基函数 $N_i$ 本身，来当质检用的权函数 $w_i$！}
    也就是令 $w_i = N_i$。
    
    \item \textbf{为什么 Galerkin 法这么牛？}
    \begin{enumerate}
    \item \textbf{对称之美：} 在力学和热传导问题中（微分算子 $\mathcal{L}$ 是自伴随/对称的），当你把 $w=N_i$ 乘进去，并结合上一章的\textbf{“分部积分大法（甩锅）”}后，你推导出来的刚度矩阵 $\bm{K}$ 一定是对称的（即 $K_{ij} = K_{ji}$）。这对计算机求解超大型矩阵极其重要！
        \item \textbf{暗合能量极值：} 对于弹性力学问题，用纯数学的 Galerkin 法得出的离散方程，恰好与物理学家用“最小势能原理”或“虚功原理”推导出的方程\textbf{一模一样}。数学与物理在 Galerkin 法中达到了完美的统一。
    \end{enumerate}
    \item \textbf{一句话总结 Galerkin 法的精髓：} 
    有限元中，要寻找的解所在的“试函数空间”，和用来检测残差的“权函数空间”，是同一个空间。
\end{itemize}

\section{第三章：变分法、虚功原理与能量观点}
如果说第二章的加权余量法是纯粹的数学“外功”，那么变分法与虚功原理就是力学学者的“内功心法”。本章要把冷冰冰的数学方程赋予“能量极值”与“做功平衡”的物理灵魂。掌握这部分，你就能深刻理解为什么力学有限元的刚度矩阵总是对称的，以及 Galerkin 法在力学中到底意味着什么。

\subsection{变分算子与泛函极值}
微积分研究的是\textbf{函数}的极值（自变量 $x$ 变了一点点，函数 $f(x)$ 怎么变）。而变分法研究的是\textbf{泛函 (Functional)} 的极值。泛函是“函数的函数”，比如一根悬链线的总长度、一个弹性系统的总能量，它的输入是一个完整的状态函数 $\bm{u}(\bm{x})$，输出是一个实数（能量）。

\textbf{1. 变分算子的物理意义：}
我们引入变分算子 $\delta$。你可以把 $\delta \bm{u}$ 想象成对真实位移场 $\bm{u}(\bm{x})$ 施加的一个极其微小的、假想的“扰动”。在力学中，我们亲切地称这个扰动为\textbf{虚位移 (Virtual Displacement)}。
\begin{itemize}
    \item \textbf{硬性规定（边界条件）：} 就像 Galerkin 法中的权函数一样，在已知真实位移的边界（本质边界条件 $\Gamma_u$）上，物体已经被死死固定住了，不允许有任何扰动。因此，在 $\Gamma_u$ 上，强制规定 $\delta \bm{u} = \bm{0}$。
    \item \textbf{运算法则：} 变分算子 $\delta$ 的运算法则和微分算子 $d$ 几乎一模一样！例如 $\delta(\bm{u} + \bm{v}) = \delta \bm{u} + \delta \bm{v}$，$\delta(\bm{u}^2) = 2\bm{u}\delta \bm{u}$。
    \item \textbf{核心铁律（算子交换）：} 变分操作与空间求导操作是可以互换顺序的！即 $\delta(\nabla \bm{u}) = \nabla (\delta \bm{u})$。翻译成力学语言就是：\textbf{扰动的应变，等于虚位移求导}（$\delta \bm{\varepsilon} = \bm{L} \delta \bm{u}$，$\bm{L}$ 为微分算子矩阵）。
\end{itemize}

\textbf{2. 泛函极值条件：}
微积分中求极值是令一阶导数 $df=0$。在变分法中，如果一个系统的能量泛函 $\Pi(\bm{u})$ 达到极值，那么对于任意容许的微小扰动 $\delta \bm{u}$，能量的一阶变分（你可以理解为泛函的“导数”）必须为 0，即：
\[ \delta \Pi = 0 \]

\subsection{最小势能原理 (Principle of Minimum Potential Energy)}
对于弹性力学系统，大自然有一个极其“吝啬”的法则：系统受力后，总是会发生变形并停留在\textbf{总势能最小}的平衡状态上。这被称为最小势能原理。

\textbf{1. 构造总势能泛函 $\Pi$：}
系统的总势能 $\Pi$ 等于积蓄在物体内部的\textbf{应变能 $U$} 减去外力对物体做的\textbf{外力势能 $W$}（外力势能等于负的外力做功）。
\[ \Pi(\bm{u}) = U(\bm{u}) - W(\bm{u}) \]
\begin{itemize}
    \item \textbf{内部应变能：} 对于线弹性体，能量密度是 $\frac{1}{2} \bm{\varepsilon}^T \bm{\sigma}$。代入本构方程 $\bm{\sigma} = \bm{D} \bm{\varepsilon}$，总应变能为：
    \[ U = \frac{1}{2} \int_{\Omega} \bm{\varepsilon}^T \bm{\sigma} \, d\Omega = \frac{1}{2} \int_{\Omega} \bm{\varepsilon}^T \bm{D} \bm{\varepsilon} \, d\Omega \]
    \item \textbf{外力势能：} 体力 $\bm{b}$ 和面力 $\bar{\bm{t}}$ 在真实位移 $\bm{u}$ 上做的功为：
    \[ W = \int_{\Omega} \bm{u}^T \bm{b} \, d\Omega + \int_{\Gamma_t} \bm{u}^T \bar{\bm{t}} \, d\Gamma \]
\end{itemize}
合并起来，弹性系统的总势能泛函为：
\[ \Pi(\bm{u}) = \frac{1}{2} \int_{\Omega} \bm{\varepsilon}^T \bm{D} \bm{\varepsilon} \, d\Omega - \int_{\Omega} \bm{u}^T \bm{b} \, d\Omega - \int_{\Gamma_t} \bm{u}^T \bar{\bm{t}} \, d\Gamma \]

\textbf{2. 求极值过程（对总势能取变分）：}
我们让系统发生一个虚位移 $\delta \bm{u}$，相应地产生虚应变 $\delta \bm{\varepsilon}$。令 $\delta \Pi = 0$：
\[ \delta \Pi = \int_{\Omega} \delta \bm{\varepsilon}^T \bm{D} \bm{\varepsilon} \, d\Omega - \int_{\Omega} \delta \bm{u}^T \bm{b} \, d\Omega - \int_{\Gamma_t} \delta \bm{u}^T \bar{\bm{t}} \, d\Gamma = 0 \]
注意到公式中的 $\bm{D} \bm{\varepsilon}$ 就是真实的应力 $\bm{\sigma}$。
如果你对这个极值方程使用第一章学过的\textbf{分部积分大法}（把 $\delta \bm{\varepsilon}$ 里的空间求导甩锅给 $\bm{\sigma}$），你会惊讶地发现，它神奇地推导出了弹性力学的强形式平衡方程 $\nabla \cdot \bm{\sigma} + \bm{b} = \bm{0}$ 以及面力边界条件 $\bm{\sigma} \cdot \bm{n} = \bar{\bm{t}}$！
\textbf{这证明了一个伟大的结论：寻找势能最小的位移场 $\bm{u}$，和求解微分平衡方程，在数学和物理上是完全等价的！}（这就是里兹法 Ritz Method 的基础）。

\subsection{虚功原理}
最小势能原理非常优雅，但它有一个致命的缺陷：\textbf{它只适用于存在“能量泛函”的保守系统}。如果你的材料进入了塑性阶段（能量被耗散成了热），或者遇到了非保守力（如随动摩擦力），势能就不存在了，最小势能原理直接罢工。

这时候，力学真正的“绝对防御”——\textbf{虚功原理}登场了。它不要求系统保守，只要系统处于平衡状态，它就永远成立，是固体力学有限元的绝对基石。

\textbf{1. 虚功原理的物理表述：}
如果一个可变形固体处于平衡状态，那么对于\textbf{任意}一个满足几何约束（在固定边界上 $\delta \bm{u} = \bm{0}$）的\textbf{虚位移场 $\delta \bm{u}$}，系统内部应力在虚应变上做的\textbf{内部虚功 (Internal Virtual Work)}，必然等于外部所有载荷在虚位移上做的\textbf{外部虚功 (External Virtual Work)}。

\textbf{2. 核心公式：}
\[ \delta W_{int} = \delta W_{ext} \]
即：
\[ \int_{\Omega} \delta \bm{\varepsilon}^T \bm{\sigma} \, d\Omega = \int_{\Omega} \delta \bm{u}^T \bm{b} \, d\Omega + \int_{\Gamma_t} \delta \bm{u}^T \bar{\bm{t}} \, d\Gamma \]
\textit{注：你看！对于线弹性系统，虚功原理的公式和前面 $\delta \Pi = 0$ 的公式长得完全一样。但虚功方程中的 $\bm{\sigma}$ 不强制要求是 $\bm{D}\bm{\varepsilon}$，它可以是任何复杂的非线性本构，这也是非线性有限元的迭代起点。}

\subsection{加权余量法与虚功原理的统一}
如果你把这一节的虚功方程，和上一节 Galerkin 法通过加权余量和分部积分得到的“弱形式方程”放在一起对比，你会发现一个震撼的事实：

\textbf{如果在数学推导中，你把“权函数 $\bm{w}$”看作力学里的“虚位移 $\delta \bm{u}$”，把“试函数 $\bm{u}$”看作“真实位移”，那么 Galerkin 法的方程与虚功原理的方程不仅形式上完全一致，连每一项的意义都严丝合缝！}
\begin{itemize}
    \item \textbf{纯数学视角：} 这是一个为了求解微分方程，用权函数 $w$ 在积分域内强迫残差平均为 0 的弱形式方程。
    \item \textbf{纯物理视角：} 这是系统处于平衡时，内力在任意虚位移 $\delta u$ 上做的功等于外力虚功的能量守恒方程。
\end{itemize}

**这就是现代有限元最迷人的地方：纯粹数学逻辑推导出的“加权余量弱形式”，在固体力学中找到了完美的物理灵魂（虚功原理）。两者殊途同归，在底层逻辑上达到了大一统。**

\section{第四章：函数空间、弱形式与连续性要求}
在完成前三章之后，我们已经知道如何从强形式走到加权积分、再走到变分与虚功。但如果要把这些推导真正变成严谨的有限元理论，还必须回答三个问题：\textbf{解究竟存在于什么函数空间中？弱形式到底把导数要求降到了什么程度？不同控制方程又会逼出什么样的连续性要求？}

\subsection{Sobolev 空间与弱导数}
为了保证能量积分收敛，我们需要明确在什么空间里去寻找有限元解。
\begin{itemize}
    \item \textbf{$L^2(\Omega)$ 平方可积空间：}
    集合中的函数 $v$ 满足 $\int_{\Omega} v^2 d\Omega < \infty$。可以理解为物理上的“有限能量”。
    \item \textbf{Sobolev 空间 $H^1(\Omega)$：}
    仅仅函数平方可积不够，有限元公式里还有梯度项 $\nabla v$。$H^1(\Omega)$ 空间要求函数本身及其一阶弱导数都在 $L^2(\Omega)$ 中：
    \[
    H^1(\Omega) = \{ v \in L^2(\Omega) \mid \nabla v \in [L^2(\Omega)]^n \}
    \]
    “弱导数”允许函数在少数点上不可导（例如连续但有转折点的折线），这完美契合了有限元中跨单元边界一阶导数不连续的特点。
    \item \textbf{$H^1_0(\Omega)$ 空间：}
    如果一个函数在 $H^1(\Omega)$ 中，并且在边界上（或 Dirichlet 边界上）等于 0，则属于 $H^1_0(\Omega)$ 空间。这正是我们选择试函数或权函数时最常出现的空间。
\end{itemize}

\subsection{弱形式的函数空间框架}
弱形式的本质，不只是“把微分方程乘个权函数再积分”，更是\textbf{把问题重新安放到一个更宽松的函数空间里}。例如：
\begin{itemize}
    \item 对泊松型二阶问题，分部积分一次后，弱形式只要求 $u$ 与 $w$ 具有一阶弱导数，因此通常工作在 $H^1(\Omega)$ 中；
    \item 对满足齐次本质边界条件的权函数，通常要求 $w \in H^1_0(\Omega)$；
    \item 若是四阶问题，如欧拉-伯努利梁和 Kirchhoff 薄板，弱形式中往往仍然保留二阶导数，因此自然会把函数空间提高到类似 $H^2$ 的层级。
\end{itemize}

\textbf{这正是“强形式 $\rightarrow$ 弱形式”最深的一层含义：}
\begin{itemize}
    \item 强形式要求逐点满足高阶微分方程；
    \item 弱形式要求在某个函数空间里满足积分恒等式；
    \item 函数空间的选择，决定了允许什么样的近似函数进入有限元离散。
\end{itemize}

\subsection{连续性要求与 $C^0/C^1$ 单元}
前面我们已经通过分部积分把强形式中的高阶导数降了阶，但“降到几阶”直接决定了离散函数至少要具备什么连续性。
\begin{itemize}
    \item 若弱形式只含一阶导数，如泊松型方程、线弹性位移型方程，则解空间通常是 $H^1(\Omega)$。这意味着离散函数只需要函数本身在单元交界处连续，而一阶导数允许跨单元跳跃。这正对应最常见的\textbf{$C^0$ 连续单元}。
    \item 线单元、三角形、四边形、四面体、六面体实体单元大多属于这一类。节点共享保证位移或温度连续，但应变、应力、热流通常在单元边界上允许不连续。
    \item 若弱形式中仍然保留二阶导数，如 Euler-Bernoulli 梁、Kirchhoff 薄板等四阶问题，则形函数不仅位移要连续，还要转角或一阶导数连续，也就是\textbf{$C^1$ 连续}。
    \item $C^1$ 形函数构造更复杂，自由度设计与单元拼接也更麻烦，因此工程中经常会采用混合形式、降阶形式或 Reissner-Mindlin 类型理论，把高连续性要求转写成 $C^0$ 可实现的形式。
\end{itemize}

\textbf{所以一定要记住：}
连续性要求不是“单元偏好”，而是由\textbf{弱形式里最高阶导数}决定的。你选什么单元，本质上是在回答：\textbf{离散空间至少要满足怎样的函数空间正则性？}

\section{第五章：有限元离散化与数值实现基础}
在完成前三章的理论铺垫，并明确第四章的函数空间与连续性要求之后，我们终于可以回答程序实现中最现实的问题：\textbf{全局矩阵和载荷向量，究竟是怎样从单元、积分点和自由度一步步“算”出来的？}

\subsection{参考单元与等参映射}
在理论推导里，积分区域是几何上真实存在的单元 $\Omega^e$；但在程序实现里，我们几乎从不直接在歪斜四边形、六面体或曲边单元上手工积分，而是统一把它们映射到一个\textbf{标准参考单元}上。例如：
\begin{itemize}
    \item 一维线单元映射到 $\xi \in [-1,1]$；
    \item 二维四边形单元映射到 $(\xi,\eta)\in[-1,1]^2$；
    \item 三维六面体单元映射到 $(\xi,\eta,\zeta)\in[-1,1]^3$。
\end{itemize}

\textbf{1. 等参元 (Isoparametric Element) 的核心思想：}
几何坐标与场变量插值使用同一套形函数 $N_i$。若单元有 $n_e$ 个节点，则真实坐标由
\[
\bm{x}(\bm{\xi}) = \sum_{i=1}^{n_e} N_i(\bm{\xi}) \bm{x}_i
\]
给出，其中 $\bm{x}_i$ 是节点的真实坐标，$\bm{\xi}$ 表示参考坐标。

\textbf{为什么这个想法如此重要？}
\begin{itemize}
    \item \textbf{统一实现：} 不管单元在真实空间里多么歪、多么扭，程序都只需在标准单元上写一套形函数与积分代码。
    \item \textbf{适配曲边几何：} 二次、三次等高阶单元可以自然描述曲边边界与复杂几何。
    \item \textbf{理论与实现一致：} 几何映射、场变量插值、导数变换三者被一套形函数体系绑在一起，不容易在代码里写乱。
\end{itemize}

\textbf{2. 雅可比矩阵与导数变换：}
从参考坐标到真实坐标的映射，关键在雅可比矩阵
\[
\bm{J} = \frac{\partial \bm{x}}{\partial \bm{\xi}}.
\]
例如在二维情形，
\[
\bm{J} =
\begin{bmatrix}
\dfrac{\partial x}{\partial \xi} & \dfrac{\partial x}{\partial \eta} \\
\dfrac{\partial y}{\partial \xi} & \dfrac{\partial y}{\partial \eta}
\end{bmatrix}.
\]
形函数在真实坐标中的梯度，并不是直接对 $x,y,z$ 求导，而是通过链式法则从参考坐标导数变换而来：
\[
\nabla_{\bm{x}} N_i = \bm{J}^{-T} \nabla_{\bm{\xi}} N_i.
\]

\textbf{3. 积分微元的变换：}
面积或体积微元满足
\[
d\Omega = \det(\bm{J}) \, d\xi d\eta d\zeta.
\]
因此代码里每个积分点乘上的 \texttt{detJ} 不是经验修正项，而是变量替换的数学结果。

\textbf{4. 工程上的危险信号：}
\begin{itemize}
    \item 若 $\det(\bm{J}) \le 0$，通常意味着单元翻转、节点编号方向错误或网格严重畸变；
    \item 若 $\det(\bm{J})$ 非常小，意味着单元接近退化，局部刚度阵容易病态；
    \item 若 $\bm{J}$ 在单元内变化剧烈，则说明几何畸变较大，即便程序不报错，精度也可能明显恶化。
\end{itemize}
所以，\textbf{检查雅可比行列式的符号与大小}，是高质量有限元程序的基本自检动作。

\subsection{单元刚度矩阵与载荷向量}
回忆最典型的弱形式结构：
\[
\int_{\Omega} \nabla w \cdot (k \nabla u)\, d\Omega
=
\int_{\Omega} w f\, d\Omega + \int_{\Gamma_q} w \bar{q}\, d\Gamma.
\]
在单元 $e$ 上做 Galerkin 离散后，取
\[
u^h|_{\Omega^e} = \bm{N}\bm{d}^e, \qquad
w^h|_{\Omega^e} = \bm{N}\delta \bm{d}^e,
\]
其中 $\bm{N}=[N_1, N_2, \dots, N_{n_e}]$，$\bm{d}^e$ 是该单元节点未知量组成的局部向量。

对于标量问题，可把梯度矩阵记作
\[
\bm{B} =
\begin{bmatrix}
\nabla N_1 & \nabla N_2 & \cdots & \nabla N_{n_e}
\end{bmatrix}.
\]
于是单元刚度阵与单元载荷向量写成统一形式：
\[
\bm{K}^e = \int_{\Omega^e} \bm{B}^T k \bm{B}\, d\Omega,
\]
\[
\bm{F}^e = \int_{\Omega^e} \bm{N}^T f\, d\Omega + \int_{\Gamma_q^e} \bm{N}^T \bar{q}\, d\Gamma.
\]

\textbf{这两个公式非常值得你对照代码理解：}
\begin{itemize}
    \item $\bm{B}$ 来自形函数求导与雅可比变换；
    \item $k$ 是材料参数或传导系数；
    \item $\bm{N}^T f$ 与 $\bm{N}^T \bar{q}$ 对应体载荷和边界载荷向节点的分配；
    \item 每一个单元都独立计算出自己的 $\bm{K}^e,\bm{F}^e$，然后等待组装。
\end{itemize}

如果问题换成线弹性力学，那么结构完全不变，只是把标量导热里的 $k$ 换成材料矩阵 $\bm{D}$，把梯度矩阵换成应变矩阵 $\bm{B}$，得到经典公式
\[
\bm{K}^e = \int_{\Omega^e} \bm{B}^T \bm{D}\bm{B}\, d\Omega.
\]

\subsection{高斯积分}
单元积分公式往往很复杂，尤其当几何映射、材料参数、形函数导数同时出现时，解析积分几乎不可行。因此有限元程序普遍采用\textbf{高斯积分 (Gaussian Quadrature)}。

\textbf{1. 一维高斯积分公式：}
\[
\int_{-1}^{1} g(\xi)\, d\xi \approx \sum_{k=1}^{n_g} w_k\, g(\xi_k),
\]
其中 $\xi_k$ 是高斯点，$w_k$ 是权重。

\textbf{2. 多维情形：}
二维和三维张量积单元常用张量积高斯积分。例如二维四边形单元中，
\[
\int_{-1}^{1}\int_{-1}^{1} g(\xi,\eta)\, d\xi d\eta
\approx
\sum_{i=1}^{n_g}\sum_{j=1}^{n_g} w_i w_j\, g(\xi_i,\eta_j).
\]
对三维六面体单元则继续扩展为三重求和。

\textbf{3. 把映射代进去之后，单元积分真正长什么样：}
\[
\bm{K}^e
=
\int_{\hat{\Omega}}
\bm{B}(\bm{\xi})^T
k
\bm{B}(\bm{\xi})
\det(\bm{J}(\bm{\xi}))
\,
d\hat{\Omega}
\approx
\sum_{q=1}^{n_q}
\bm{B}_q^T k \bm{B}_q \det(\bm{J}_q) w_q.
\]
这正是代码里最常见的积分点循环来源。

\textbf{4. 为什么积分点数不能随便选？}
\begin{itemize}
    \item \textbf{点太少：} 会欠积分，刚度被算软，甚至出现零能模态与沙漏模式；
    \item \textbf{点太多：} 会提高计算成本，有时对低阶单元收益有限；
    \item \textbf{点数恰当：} 才能在精度、稳定性与计算量之间取得平衡。
\end{itemize}

\textbf{5. 完全积分与减缩积分：}
\begin{itemize}
    \item \textbf{完全积分 (Full Integration)：} 使用足够多的积分点，确保给定多项式阶次下的积分精度；
    \item \textbf{减缩积分 (Reduced Integration)：} 故意少用一些积分点，以降低闭锁或计算成本；
    \item \textbf{风险：} 减缩积分虽然常能缓解闭锁，但也可能引入非物理零能模态，需要配合稳定化技术。
\end{itemize}

\subsection{总体组装与稀疏结构}
有限元程序最迷人的地方之一在于：全局大矩阵不是直接推导出来的，而是由无数个局部小矩阵按连接关系拼起来的。

\textbf{1. 节点连接表是组装的桥梁：}
设单元 $e$ 的局部节点编号为 $a=1,2,\dots,n_e$，它们在全局系统中的自由度编号由连接映射 $I^e(a)$ 给出。则组装规则是
\[
K_{I^e(a)\, I^e(b)} \mathrel{+}= K^e_{ab},
\qquad
F_{I^e(a)} \mathrel{+}= F^e_a.
\]

\textbf{2. 为什么全局矩阵天然稀疏？}
因为每个形函数只在少数相邻单元中非零，任何一个节点只会与其局部邻域内的节点发生耦合。于是：
\begin{itemize}
    \item 绝大多数 $K_{ij}$ 实际上恒为 0；
    \item 稀疏存储格式（CSR、CSC、Skyline 等）成为大规模计算的必需品；
    \item 有限元的局部支集数学性质，直接决定了程序的数据结构设计。
\end{itemize}

\textbf{3. 一个典型的单元循环长什么样：}
\begin{enumerate}
    \item 读取单元节点坐标与自由度编号；
    \item 在参考单元上遍历高斯积分点；
    \item 计算形函数、形函数导数、雅可比矩阵、$\det(\bm{J})$ 与 $\bm{B}$；
    \item 累加得到 $\bm{K}^e$ 与 $\bm{F}^e$；
    \item 用连接表把 $\bm{K}^e,\bm{F}^e$ 组装进全局系统；
    \item 处理边界条件后，调用线性方程组求解器。
\end{enumerate}

\subsection{边界条件的离散处理}
连续问题中的边界条件，到了离散系统里必须变成非常具体的矩阵操作。这里最容易“只会写代码，不知其然”，所以必须把逻辑理清。

\textbf{1. 自然边界条件：}
像 Neumann 边界、表面力、热流密度这类条件，本来就是弱形式的一部分，因此它们通常表现为右端载荷向量中的边界积分项：
\[
\bm{F}_{\Gamma}^e = \int_{\Gamma_q^e} \bm{N}^T \bar{q}\, d\Gamma.
\]
也就是说，\textbf{自然边界条件往往不需要强行施加}，因为它已经在变分形式里自然出现了。

\textbf{2. 本质边界条件：}
Dirichlet 边界、固定位移这类条件则不同，它们规定的是未知量本身的取值，必须直接进入自由度约束。常见处理方式有：
\begin{itemize}
    \item \textbf{直接消元法：} 删除已知自由度对应的方程或把其贡献移到右端；
    \item \textbf{矩阵修改法：} 将对应行列清零、对角置一、右端置为已知值；
    \item \textbf{罚函数法：} 在对角线上加入一个很大的系数强行逼近约束。
\end{itemize}

\textbf{3. 你需要建立的直觉：}
\begin{itemize}
    \item 自然边界条件是从弱形式里长出来的；
    \item 本质边界条件是对离散自由度施加约束；
    \item 这一区别不是术语游戏，而是决定了程序应该改右端还是改矩阵。
\end{itemize}

\subsection{闭锁与数值病态}
很多初学者一旦发现结果“偏硬”“位移过小”，第一反应是怀疑材料参数、边界条件或程序 bug。其实有时程序完全没错，问题出在\textbf{离散空间本身不适合该物理状态}，这就是闭锁 (Locking)。

\textbf{1. 体积闭锁：}
当材料接近不可压缩，例如泊松比 $\nu \to 0.5$ 时，低阶位移单元很难同时满足“几何兼容”和“近似零体积应变”两个要求，结果就是单元表现得异常僵硬，位移被严重低估。

\textbf{2. 剪切闭锁：}
在薄梁、薄板、薄壳问题中，若单元插值无法正确表达“薄结构中横向剪切应变应趋近于零”的物理特征，就会人为引入虚假的剪切能量，导致结构算得过硬。

\textbf{3. 常见缓解手段：}
\begin{itemize}
    \item 减缩积分；
    \item 选择性减缩积分；
    \item $\bar{B}$ 方法；
    \item 混合有限元或引入附加场变量；
    \item 选用更适合该物理问题的高阶单元。
\end{itemize}

\textbf{4. 但“解药”也可能带来副作用：}
例如过度减缩积分可能触发沙漏模式；混合单元若不满足稳定性条件，后面第六章会讲到的 LBB 条件就会出问题。也就是说，\textbf{闭锁、积分方案、稳定性}三者从来不是分开的主题，而是连续问题、离散空间与数值算法之间的联动结果。

\textbf{5. 除了闭锁，还要警惕网格质量导致的病态：}
\begin{itemize}
    \item 高长宽比单元会恶化条件数；
    \item 严重扭曲单元会破坏插值质量；
    \item 过粗网格会带来离散误差，过细网格则提高计算代价并放大求解器压力。
\end{itemize}

\textbf{第五章练习建议：}
\begin{enumerate}
    \item 亲手推导四节点双线性四边形单元的形函数、雅可比矩阵与 $\bm{B}$ 矩阵，并写出 $2\times2$ 高斯积分下的局部刚度阵表达式。
    \item 对同一个薄结构或近不可压缩问题，比较完全积分与减缩积分的计算结果，观察闭锁与过软化的差别。
    \item 对照你已有的有限元代码，逐行标记每一段循环分别对应“形函数插值”“雅可比变换”“积分点累加”“局部到全局组装”“边界条件处理”中的哪一步。
\end{enumerate}

\section{第六章：收敛性、稳定性与误差分析}
前五章解决的是“怎样把连续问题写成有限元可以求的离散系统”，而本章要回答的是另一个更尖锐的问题：\textbf{离散系统算出来的结果为什么可信？它什么时候收敛？什么时候稳定？误差又如何估计？}

\subsection{Lax-Milgram 定理}
Lax-Milgram 定理给出了弱形式存在唯一解的数学保障。若双线性型 $a(u,v)$ 满足：
\begin{itemize}
    \item \textbf{连续性：} $|a(u,v)| \le M \|u\| \|v\|$；
    \item \textbf{强制性：} $a(u,u) \ge \alpha \|u\|^2$，
\end{itemize}
那么弱形式问题就存在唯一解。它相当于告诉你：在函数空间里，这个问题没有“解不存在”或“解不唯一”的根本性灾难。

\subsection{先验误差估计}
有限元误差分析最经典的结果之一，是先验误差界：
\[
\| u - u^h \|_{H^1} \le C h^p \| u \|_{H^{p+1}}.
\]
它揭示了两件事：
\begin{itemize}
    \item \textbf{$h$ 细化：} 网格越细，误差越小；
    \item \textbf{$p$ 升阶：} 形函数阶次越高，若真解足够光滑，则收敛速度越快。
\end{itemize}
这就是有限元里著名的 $h$-version 与 $p$-version 收敛思想。

\subsection{LBB 条件}
对于混合有限元问题，例如不可压缩流体、完全不可压缩弹性等，仅仅有双线性型还不够，还需要满足 LBB 条件（又叫 Inf-Sup 条件）：
\[
\inf_{q} \sup_{\bm{v}} \frac{b(\bm{v}, q)}{\|\bm{v}\| \|q\|} \ge \beta > 0.
\]
它是混合格式稳定的充分条件。若不满足这个条件，就可能出现压力振荡、伪模态或完全错误的耦合结果。

\section{第七章：完整实战（一）—— 泊松方程}
泊松方程是传热学、静电场以及渗流问题中最经典的方程，形式上与弹性力学方程高度相似。这里我们将运用前三章的知识（散度定理、分部积分、加权余量），事无巨细地推导有限元求解的完整步骤。

\subsection{问题定义与边界条件}
假设我们要求解区域 $\Omega$ 内的稳态热传导问题。
\begin{itemize}
    \item \textbf{控制方程（偏微分方程）：}
    \[ -\nabla \cdot (k \nabla u) = f \quad \text{在 } \Omega \text{ 内} \]
    其中：$u$ 是待求的温度场；$k$ 是热导率（如果是各向同性材料就是一个常数）；$f$ 是内部热源生成率。因为方程要求 $u$ 在域内逐点满足二阶可导，所以这叫\textbf{强形式 (Strong Form)}。
    
    \item \textbf{边界条件 (Boundary Conditions, BCs)：}
    区域的边界 $\partial\Omega$ 由两部分组成：$\partial\Omega = \Gamma_u \cup \Gamma_q$，且两部分互不重叠。
    \begin{enumerate}
        \item \textbf{本质边界条件 (Dirichlet BC / Essential BC)：} 
        直接规定了边界上的温度值。
        \[ u = \bar{u} \quad \text{在 } \Gamma_u \text{ 上} \]
        在有限元里，满足这个条件的节点位移是直接已知的，不需要通过解矩阵方程求。
        
        \item \textbf{自然边界条件 (Neumann BC / Natural BC)：}
        规定了边界上的热流密度（即温度的法向导数乘上热导率）。
        \[ (k \nabla u) \cdot \bm{n} = \bar{q} \quad \text{在 } \Gamma_q \text{ 上} \]
        这里的 $\bm{n}$ 是边界单位外法向，$\bar{q}$ 是给定的热通量。自然边界条件会在弱形式的分部积分里自然出现，因此得名。
    \end{enumerate}
\end{itemize}

\subsection{弱形式构造}
既然强形式要求太高（要在无数个点上精确成立，还要算二阶导），我们要把它弱化。

\textbf{（1）引入权函数与全域积分}
设 $w$ 是一个任意的、平滑的\textbf{权函数（试函数）}。为了后续处理边界条件时方便，我们强制规定：\textbf{在已知温度的边界 $\Gamma_u$ 上，$w = 0$}。
将控制方程移项得到残差，然后在域 $\Omega$ 上用 $w$ 对残差求加权平均，并令其等于0：
\[ \int_{\Omega} w \Big( -\nabla \cdot (k \nabla u) - f \Big) \, d\Omega = 0 \]
展开并移项：
\[ -\int_{\Omega} w \nabla \cdot (k \nabla u) \, d\Omega = \int_{\Omega} w f \, d\Omega \]
此时左边还是包含 $u$ 的二阶导数（包含在 $\nabla \cdot (\dots)$ 里面）。

\textbf{（2）施展“降维魔法”（分部积分）}
利用我们在第1.3节学到的连招（散度定理+乘积求导），对左边的积分项进行“甩锅”：
\[ -\int_{\Omega} w \nabla \cdot (k \nabla u) \, d\Omega = \int_{\Omega} (\nabla w) \cdot (k \nabla u) \, d\Omega - \int_{\partial\Omega} w (k \nabla u) \cdot \bm{n} \, d\Gamma \]
将其代回之前的方程，我们得到：
\[ \int_{\Omega} (\nabla w) \cdot (k \nabla u) \, d\Omega - \int_{\partial\Omega} w (k \nabla u) \cdot \bm{n} \, d\Gamma = \int_{\Omega} w f \, d\Omega \]

\textbf{（3）代入边界条件，得到最终的弱形式}
看等号左边产生的那个边界积分项 $\int_{\partial\Omega}$，边界由两部分组成：$\Gamma_u$ 和 $\Gamma_q$。
\begin{itemize}
    \item \textbf{在 $\Gamma_u$ 上：} 我们前面规定了 $w=0$，所以这一部分的积分直接等于 $0$！（看，本质边界条件通过权函数的巧妙设置被处理掉了）。
    \item \textbf{在 $\Gamma_q$ 上：} 我们已知自然边界条件 $(k \nabla u) \cdot \bm{n} = \bar{q}$，直接代入！
\end{itemize}
于是，边界项变成了：$-\int_{\Gamma_q} w \bar{q} \, d\Gamma$。移项到右边，最终的\textbf{弱形式方程}诞生：
\[ \int_{\Omega} \nabla w \cdot (k \nabla u) \, d\Omega = \int_{\Omega} w f \, d\Omega + \int_{\Gamma_q} w \bar{q} \, d\Gamma \]
\textit{意义：} 这个积分方程里，$u$ 和 $w$ 都只有一阶导数（$\nabla u$ 和 $\nabla w$），并且包含了所有的自然边界条件。这正是物理学中“虚功原理”或“能量守恒”的数学体现！

\subsection{Galerkin 离散}
数学推导完毕，开始引入代码思维。无穷维的函数空间在计算机里算不了，必须将连续的域划分为离散的单元（网格划分），用节点上的值来近似。

\textbf{（1）试函数 $u$ 的离散（插值）：}
把未知的温度场 $u(\bm{x})$ 用节点温度 $u_j$ 和已知的形函数（基函数）$N_j(\bm{x})$ 展开：
\[ u(\bm{x}) \approx \sum_{j=1}^{n} N_j(\bm{x}) u_j \]
对其求梯度（由于节点温度 $u_j$ 只是个标量数字，空间导数只作用在形函数上）：
\[ \nabla u = \sum_{j=1}^{n} \nabla N_j u_j \]

\textbf{（2）权函数 $w$ 的选取（Galerkin 法核心）：}
Galerkin 方法规定，权函数 $w$ 的选择应该与试函数的基函数长得一模一样！也就是取 $w(\bm{x}) = N_i(\bm{x})$（其中 $i$ 代表我们正在构造第 $i$ 个方程，对应第 $i$ 个节点）。
那么它的梯度就是 $\nabla w = \nabla N_i$。

\textbf{（3）代入弱形式：}
把上面离散后的 $u$ 和 $w$ 代入弱形式方程的左边：
\[ \int_{\Omega} \nabla N_i \cdot \left( k \sum_{j=1}^{n} \nabla N_j u_j \right) \, d\Omega = \sum_{j=1}^{n} \left( \int_{\Omega} \nabla N_i \cdot (k \nabla N_j) \, d\Omega \right) u_j \]
右边（载荷项）也对应替换 $w$ 为 $N_i$：
\[ \int_{\Omega} N_i f \, d\Omega + \int_{\Gamma_q} N_i \bar{q} \, d\Gamma \]

\subsection{矩阵系统组装}
将第三步骤中对于每一个节点 $i$ ($i = 1, 2, \dots, n$) 产生的代数方程组合起来，就形成了一个庞大的线性代数方程组 $\bm{K}\bm{u} = \bm{F}$：

\begin{itemize}
    \item \textbf{刚度矩阵 $\bm{K}$（也叫传导矩阵）：}
    由于左边求和里的积分部分只跟形函数的梯度和材料参数 $k$ 有关，它是一个已知的常数，这就是刚度矩阵的第 $(i, j)$ 个元素：
    \[ K_{ij} = \int_{\Omega} \nabla N_i \cdot (k \nabla N_j) \, d\Omega \]
    \textit{代码提示：注意看公式，因为点乘满足交换律 $\nabla N_i \cdot \nabla N_j = \nabla N_j \cdot \nabla N_i$，所以 $K_{ij} = K_{ji}$。这从数学上证明了泊松方程对应的有限元刚度阵必定是一个\textbf{对称矩阵}！这在代码中能节省一半的存储空间和大量的求解时间。}
    
    \item \textbf{载荷向量 $\bm{F}$（右端项）：}
    第 $i$ 个节点上的等效载荷为：
    \[ F_i = \int_{\Omega} N_i f \, d\Omega + \int_{\Gamma_q} N_i \bar{q} \, d\Gamma \]
    第一项是体热源 $f$ 根据形函数 $N_i$ 分配到节点 $i$ 上的热量，第二项是边界热通量 $\bar{q}$ 分配到节点 $i$ 上的热量。
\end{itemize}

\textbf{最终的闭环：} 
从这一刻起，复杂的偏微分方程变成了一个纯粹的线性代数问题 $\bm{K}\bm{u} = \bm{F}$！
在实际写代码（C++ / Python / Fortran）时：
你只需要在一层 `for` 循环中遍历每一个单元 (element)，利用高斯积分算出手头上这个小单元的局部刚度阵 $\bm{K}^e$ 和局部载荷 $\bm{F}^e$，然后把它们叠加（组装）到全局的大矩阵 $\bm{K}$ 和大向量 $\bm{F}$ 中相应的位置。最后调用线性求解器解出节点温度向量 $\bm{u}$，大功告成！

\section{第八章：完整实战（二）—— 欧拉-伯努利梁}
如果说泊松方程是“二阶偏微分方程如何进入有限元”的标准教科书范例，那么欧拉-伯努利梁就是另一个极其关键的范例：\textbf{它是一个四阶问题}。这意味着你不能只做一次分部积分就收工，也不能再随便使用普通的 $C^0$ 线性插值函数。这个例子会直接把你带到有限元理论里一个非常本质的问题上：\textbf{弱形式到底把连续性要求降到了什么程度？}

下面我们用最经典的一维受弯梁问题，把“强形式 $\rightarrow$ 两次分部积分 $\rightarrow$ 弱形式 $\rightarrow$ Hermite 插值 $\rightarrow$ 单元刚度阵 $\rightarrow$ 载荷向量 $\rightarrow$ 求解”完整走一遍。

\textbf{符号约定：}
本节统一取梁轴方向为 $x\in[0,L]$，挠度为 $w(x)$，分布载荷为 $q(x)$，并约定 $q$ 的正方向与 $w$ 的正方向一致。梁的弯曲刚度为 $EI$，其中 $E$ 是弹性模量，$I$ 是截面二次矩。为避免不同教材之间正负号定义互相打架，我们直接采用最稳妥的控制方程写法：
\[
(EI\, w'')'' = q(x).
\]
这样后面的推导最清楚，也最不容易在边界项的正负号上绕晕。

\subsection{问题定义与边界条件}
欧拉-伯努利梁理论建立在一个核心假设上：\textbf{平截面在弯曲后仍保持平面且垂直于中性轴}。这意味着梁的横向剪切变形被忽略，因此弯曲能量只与曲率 $w''$ 有关。

\textbf{1. 强形式：}
\[
(EI\, w'')'' = q \qquad \text{在 } (0,L) \text{ 内}.
\]
若 $EI$ 为常数，则方程可写成
\[
EI\, w^{(4)} = q.
\]
这就是一个典型的四阶常微分方程，所以它需要四个边界条件才能唯一确定解。

\textbf{2. 梁问题中的四类边界量：}
\begin{itemize}
    \item \textbf{位移（挠度）：} $w$；
    \item \textbf{转角（斜率）：} $\theta = w'$；
    \item \textbf{弯矩：} $M = EI\, w''$；
    \item \textbf{剪力：} $V = (EI\, w'')'$。
\end{itemize}

\textbf{3. 本质边界条件与自然边界条件：}
对于梁问题，\textbf{本质边界条件}不再只有一个，而是两类：
\[
w = \bar{w}, \qquad w' = \bar{\theta}.
\]
对应地，\textbf{自然边界条件}是：
\[
M = \overline{M}, \qquad V = \overline{V}.
\]

\textbf{4. 为什么这是四阶问题的关键区别？}
在泊松方程里，位移型变量只有一个本质边界条件；而在梁里，由于控制方程是四阶，\textbf{位移和转角都可能被指定}。这一点非常重要，因为它会直接决定试函数和权函数必须满足什么样的边界约束。

\textbf{5. 典型边界条件举例：}
\begin{itemize}
    \item \textbf{固支端 (Clamped End)：} $w=0,\; w'=0$；
    \item \textbf{简支端 (Simply Supported End)：} $w=0,\; M=EI w''=0$；
    \item \textbf{自由端 (Free End)：} $M=0,\; V=0$；
    \item \textbf{端部受力/受矩：} $V=\overline{V},\; M=\overline{M}$。
\end{itemize}

\subsection{加权余量方程}
现在照着第二章加权余量法的套路来。令 $v(x)$ 为任意权函数（你也可以把它理解为虚位移函数），并要求它满足与本质边界条件相容的约束。

将强形式乘以 $v$ 并在区间 $(0,L)$ 上积分：
\[
\int_0^L v (EI\, w'')''\, dx = \int_0^L v q\, dx.
\]
如果你此时停下不动，那么左边仍然含有 $w$ 的四阶导数。计算机当然不喜欢，有限元形函数更不喜欢。我们必须把导数阶次降下来。

\subsection{第一次分部积分}
对左边做一次分部积分：
\[
\int_0^L v (EI\, w'')''\, dx
=
\left[ v (EI\, w'')' \right]_0^L
-
\int_0^L v' (EI\, w'')'\, dx.
\]
你现在会发现：$w$ 的四阶导数降成了三阶导数，但这还不够。因为积分项中仍然有 $(EI\,w'')'$，也就是三阶导数量级。对于标准位移型有限元来说，这个要求依然太高。

\subsection{第二次分部积分}
对剩余积分项再做一次分部积分：
\[
-\int_0^L v' (EI\, w'')'\, dx
=
-\left[ v' EI\, w'' \right]_0^L
+
\int_0^L v'' EI\, w''\, dx.
\]
代回去，得到
\[
\int_0^L v (EI\, w'')''\, dx
=
\left[ v (EI\, w'')' - v' EI\, w'' \right]_0^L
+
\int_0^L EI\, v'' w''\, dx.
\]
再与右端载荷项相等，可得
\[
\int_0^L EI\, v'' w''\, dx
=
\int_0^L v q\, dx
-
\left[ v (EI\, w'')' - v' EI\, w'' \right]_0^L.
\]

\textbf{这一步是整个梁有限元推导的灵魂：}
\begin{itemize}
    \item 强形式里是四阶导数 $w^{(4)}$；
    \item 弱形式里只剩二阶导数 $w''$ 与 $v''$；
    \item 边界上自然冒出来两个边界项，一个乘 $v$，一个乘 $v'$；
    \item 这正对应了梁问题有两类自然边界量：剪力和弯矩。
\end{itemize}

\subsection{弱形式与自然边界条件}
把边界项写成更有物理意义的样子。记
\[
M = EI\, w'', \qquad V = (EI\, w'')'.
\]
则上式变成
\[
\int_0^L EI\, v'' w''\, dx
=
\int_0^L v q\, dx
-
\left[ v V - v' M \right]_0^L.
\]

如果你考虑一根\textbf{左端固支、右端可能受端力与端矩}的悬臂梁，那么在 $x=0$ 处有
\[
w(0)=0,\qquad w'(0)=0.
\]
因为权函数必须与本质边界条件相容，所以
\[
v(0)=0,\qquad v'(0)=0.
\]
此时边界项只剩右端 $x=L$ 的贡献：
\[
\int_0^L EI\, v'' w''\, dx
=
\int_0^L v q\, dx
- v(L)\,V(L) + v'(L)\,M(L).
\]

若右端施加已知端剪力 $\overline{V}_L$ 和端弯矩 $\overline{M}_L$，则可写成
\[
\int_0^L EI\, v'' w''\, dx
=
\int_0^L v q\, dx
- v(L)\,\overline{V}_L + v'(L)\,\overline{M}_L.
\]

\textbf{这告诉你一件很深刻的事：}
\begin{itemize}
    \item 与 $v$ 相乘的边界量是剪力；
    \item 与 $v'$ 相乘的边界量是弯矩；
    \item 所以在梁问题里，权函数本身和它的一阶导数都会出现在边界虚功项中。
\end{itemize}

\textbf{一句话总结弱形式：}
寻找满足本质边界条件的 $w$，使得对所有满足齐次本质边界条件的权函数 $v$，都有
\[
\int_0^L EI\, v'' w''\, dx
=
\int_0^L v q\, dx
- v(L)\,\overline{V}_L + v'(L)\,\overline{M}_L.
\]
这就是欧拉-伯努利梁的标准弱形式。

\subsection{连续性要求与 $C^1$ 条件}
观察弱形式左端：
\[
\int_0^L EI\, v'' w''\, dx.
\]
这里要求试函数 $w$ 至少具有平方可积的二阶导数，也就是说从函数空间角度看，它应该属于类似 $H^2$ 的空间，而不是泊松方程里常见的 $H^1$。

这意味着：
\begin{itemize}
    \item $w$ 本身在单元交界处要连续；
    \item $w'$ 在单元交界处也要连续；
    \item 因而不能只使用普通的 $C^0$ 拉格朗日线单元；
    \item 必须使用至少 $C^1$ 连续的插值函数。
\end{itemize}

\textbf{这就是为什么梁/板的经典位移型有限元，经常会引出 Hermite 形函数。}
泊松方程教会你“弱形式把二阶问题降到一阶导”；欧拉-伯努利梁进一步告诉你：“即使降完阶，四阶问题依旧要求更高的连续性。”

\subsection{Hermite 梁单元插值}
取一个长度为 $L_e$ 的梁单元，局部坐标为 $x\in[0,L_e]$。与普通线单元不同，梁单元的每个节点有两个自由度：
\[
\{d^e\} =
\begin{bmatrix}
w_1 & \theta_1 & w_2 & \theta_2
\end{bmatrix}^T,
\]
其中 $w_1,w_2$ 是节点挠度，$\theta_1,\theta_2$ 是节点转角。

令无量纲坐标
\[
\xi = \frac{x}{L_e}, \qquad 0 \le \xi \le 1.
\]
三次 Hermite 形函数写成
\[
w^h(x) = N_1(\xi) w_1 + N_2(\xi) \theta_1 + N_3(\xi) w_2 + N_4(\xi) \theta_2.
\]
其中
\[
N_1 = 1 - 3\xi^2 + 2\xi^3,
\]
\[
N_2 = L_e(\xi - 2\xi^2 + \xi^3),
\]
\[
N_3 = 3\xi^2 - 2\xi^3,
\]
\[
N_4 = L_e(-\xi^2 + \xi^3).
\]

\textbf{为什么第二和第四个形函数前面要乘 $L_e$？}
因为它们对应的自由度不是位移，而是转角 $\theta = dw/dx$，它本身量纲与位移不同。乘上 $L_e$ 之后，整个插值式在量纲上才一致。

\textbf{这些 Hermite 形函数有四个关键插值性质：}
\begin{itemize}
    \item 在节点 1 处，$N_1=1$ 且其余对位移自由度的贡献为 0；
    \item 在节点 1 处，$N_2$ 对应单位转角插值；
    \item 在节点 2 处，$N_3=1$；
    \item 在节点 2 处，$N_4$ 对应单位转角插值。
\end{itemize}
它们不仅保证挠度连续，还保证转角连续，因此恰好满足梁弱形式所需的 $C^1$ 连续性。

\subsection{曲率矩阵构造}
欧拉-伯努利梁的应变能与曲率 $\kappa = w''$ 有关，因此我们要对插值函数求二阶导数：
\[
w^h = \bm{N}\bm{d}^e,
\qquad
w^{h\,\prime\prime} = \bm{B}\bm{d}^e.
\]
其中
\[
\bm{N} =
\begin{bmatrix}
N_1 & N_2 & N_3 & N_4
\end{bmatrix},
\qquad
\bm{B} =
\begin{bmatrix}
N_1'' & N_2'' & N_3'' & N_4''
\end{bmatrix}.
\]
由于 $\xi = x/L_e$，所以
\[
\frac{d}{dx} = \frac{1}{L_e}\frac{d}{d\xi},
\qquad
\frac{d^2}{dx^2} = \frac{1}{L_e^2}\frac{d^2}{d\xi^2}.
\]
把四个形函数逐个求二阶导，可得
\[
N_1'' = \frac{-6 + 12\xi}{L_e^2},
\qquad
N_2'' = \frac{-4 + 6\xi}{L_e},
\]
\[
N_3'' = \frac{6 - 12\xi}{L_e^2},
\qquad
N_4'' = \frac{-2 + 6\xi}{L_e}.
\]
因此
\[
\bm{B} =
\begin{bmatrix}
\dfrac{-6 + 12\xi}{L_e^2} &
\dfrac{-4 + 6\xi}{L_e} &
\dfrac{6 - 12\xi}{L_e^2} &
\dfrac{-2 + 6\xi}{L_e}
\end{bmatrix}.
\]

\subsection{单元刚度矩阵}
将离散形式代入弱形式左端，单元弯曲刚度矩阵为
\[
\bm{K}^e = \int_0^{L_e} EI\, \bm{B}^T \bm{B}\, dx.
\]
若 $EI$ 在单元内为常数，经过积分可得到经典的两节点 Euler-Bernoulli 梁单元刚度阵：
\[
\bm{K}^e
=
\frac{EI}{L_e^3}
\begin{bmatrix}
12 & 6L_e & -12 & 6L_e \\
6L_e & 4L_e^2 & -6L_e & 2L_e^2 \\
-12 & -6L_e & 12 & -6L_e \\
6L_e & 2L_e^2 & -6L_e & 4L_e^2
\end{bmatrix}.
\]

\textbf{这个矩阵每一项都很有物理味道：}
\begin{itemize}
    \item 对角项代表节点本身挠度/转角对内力的刚性贡献；
    \item 非对角项代表节点之间的弯曲耦合；
    \item 矩阵对称，源于虚功原理与双线性型的对称性；
    \item 整体比例是 $EI/L_e^3$，这非常符合直觉：梁越硬（$EI$ 越大）越难弯，梁越短（$L_e$ 越小）刚度急剧增大。
\end{itemize}

\subsection{一致载荷向量}
弱形式右端的分布载荷项为
\[
\int_0^{L_e} \bm{N}^T q(x)\, dx.
\]
因此单元一致载荷向量定义为
\[
\bm{F}^e_q = \int_0^{L_e} \bm{N}^T q(x)\, dx.
\]
若单元上承受常值均布载荷 $q(x)=q_0$，积分可得
\[
\bm{F}^e_q
=
q_0
\begin{bmatrix}
\dfrac{L_e}{2} \\
\dfrac{L_e^2}{12} \\
\dfrac{L_e}{2} \\
-\dfrac{L_e^2}{12}
\end{bmatrix}.
\]

\textbf{注意这里最容易犯的错误：}
\begin{itemize}
    \item 很多人只把均布载荷“平均分给两个节点”，写成 $\left[q_0L_e/2,\;0,\;q_0L_e/2,\;0\right]^T$；
    \item 这对梁单元是不对的，因为端部转角自由度也会通过虚功参与做功；
    \item 正确的一致载荷向量必须由 $\int \bm{N}^T q\, dx$ 推导出来，所以会自然出现节点力矩项 $\pm q_0L_e^2/12$。
\end{itemize}

如果右端还作用集中端剪力 $\overline{V}_L$ 与端弯矩 $\overline{M}_L$，则它们在单元右端自由度上的贡献可直接写入等效节点载荷：
\[
\bm{F}^e_{\text{end}} =
\begin{bmatrix}
0 \\ 0 \\ -\overline{V}_L \\ \overline{M}_L
\end{bmatrix},
\]
其中正负号必须与你采用的力和位移正方向约定保持一致。真正写程序时，\textbf{最稳妥的做法永远是回到弱形式边界项逐项核对}，而不是凭印象背符号。

\subsection{全局组装与边界条件}
到这里，梁问题已经完全进入有限元标准流程。对每个梁单元，你都可以算出
\[
\bm{K}^e, \qquad \bm{F}^e = \bm{F}^e_q + \bm{F}^e_{\text{end}}.
\]
然后根据节点连接表把它们组装到全局系统：
\[
\bm{K}\bm{d} = \bm{F}.
\]
这里的全局未知向量不是单纯的位移，而是
\[
\bm{d} =
\begin{bmatrix}
w_1 & \theta_1 & w_2 & \theta_2 & \cdots
\end{bmatrix}^T.
\]
所以梁单元程序在自由度编号时，通常采用“每节点两个自由度”的布局方式。

本质边界条件的施加方式也与泊松方程类似，只不过现在要处理两类自由度：
\begin{itemize}
    \item 若端部固支，则同时约束该节点的 $w$ 与 $\theta$；
    \item 若端部简支，则约束 $w$，但转角 $\theta$ 不约束；
    \item 若端部自由，则不约束位移型自由度，而是通过自然边界项体现外载。
\end{itemize}

\subsection{悬臂梁算例}
考虑一根长度为 $L$、弯曲刚度为常数 $EI$ 的悬臂梁。左端固支，右端自由，全梁承受均布载荷 $q_0$。我们先用\textbf{一个两节点 Hermite 梁单元}离散整个梁，看看有限元方程如何落地。

\textbf{1. 单元刚度阵：}
\[
\bm{K}^e
=
\frac{EI}{L^3}
\begin{bmatrix}
12 & 6L & -12 & 6L \\
6L & 4L^2 & -6L & 2L^2 \\
-12 & -6L & 12 & -6L \\
6L & 2L^2 & -6L & 4L^2
\end{bmatrix}.
\]

\textbf{2. 单元一致载荷向量：}
\[
\bm{F}^e
=
q_0
\begin{bmatrix}
\dfrac{L}{2} \\
\dfrac{L^2}{12} \\
\dfrac{L}{2} \\
-\dfrac{L^2}{12}
\end{bmatrix}.
\]

\textbf{3. 施加左端固支条件：}
左端节点 1 有
\[
w_1 = 0,\qquad \theta_1 = 0.
\]
因此未知量只剩右端自由度
\[
\begin{bmatrix}
w_2 \\ \theta_2
\end{bmatrix}.
\]
约化后的线性系统为
\[
\frac{EI}{L^3}
\begin{bmatrix}
12 & -6L \\
-6L & 4L^2
\end{bmatrix}
\begin{bmatrix}
w_2 \\ \theta_2
\end{bmatrix}
=
q_0
\begin{bmatrix}
\dfrac{L}{2} \\
-\dfrac{L^2}{12}
\end{bmatrix}.
\]

\textbf{4. 求解得到自由端位移与转角：}
解这个 $2\times2$ 线性方程组，可得
\[
w_2 = \frac{q_0 L^4}{8EI},
\qquad
\theta_2 = \frac{q_0 L^3}{6EI}.
\]

\textbf{5. 这两个结果意味着什么？}
\begin{itemize}
    \item 它们正好等于悬臂梁在均布载荷下的经典解析解的端部挠度与端部转角；
    \item 这说明两节点 Hermite 梁单元不仅形式漂亮，而且对很多基础梁问题非常有效；
    \item 但要注意：即便端部位移算得很准，单元内部弯矩/剪力分布仍然是有限元插值意义下的近似，需要通过后处理恢复。
\end{itemize}

\subsection{弯矩与剪力后处理}
有限元主求解阶段给我们的直接结果，其实只是全局自由度向量 $\bm{d}$，也就是各节点的挠度和转角。工程上真正关心的，往往却是：
\begin{itemize}
    \item 梁哪里弯矩最大；
    \item 哪个截面剪力最大；
    \item 哪一段最容易先屈服或开裂。
\end{itemize}
所以求出位移以后，必须再做一步\textbf{后处理恢复 (Post-processing Recovery)}，把位移型结果变成内力型结果。

\textbf{这一节最核心的一句话是：}
\[
\bm{d} \quad \Longrightarrow \quad w(x),\; w'(x),\; w''(x),\; M(x),\; V(x).
\]
也就是说，后处理本质上就是“对离散位移场做导数与物理量映射”。

\textbf{1. 从位移插值得到曲率：}
对任一单元，
\[
w^h(x) = \bm{N}(x)\bm{d}^e.
\]
因此
\[
\theta^h(x)=w^{h\,\prime}(x),
\qquad
\kappa^h(x)=w^{h\,\prime\prime}(x)=\bm{B}(x)\bm{d}^e.
\]
这里 $\kappa=w''$ 就是梁的曲率。

\textbf{2. 由曲率恢复弯矩：}
根据欧拉-伯努利梁本构关系，
\[
M(x)=EI\,\kappa(x)=EI\,w''(x).
\]
所以单元内的有限元弯矩场可直接写成
\[
M^h(x)=EI\,\bm{B}(x)\bm{d}^e.
\]
这一步非常直接：\textbf{只要知道单元自由度向量 $\bm{d}^e$，在任意一个 $x$ 点都能算出弯矩。}

\textbf{3. 由弯矩继续恢复剪力：}
按照本节前面的符号约定，
\[
V(x) = (EI\, w'')' = M'(x).
\]
若 $EI$ 为常数，则
\[
V^h(x)=EI\,w^{h\,\prime\prime\prime}(x).
\]
进一步写成矩阵形式就是
\[
V^h(x)=EI\,\bm{B}_{,x}(x)\bm{d}^e,
\]
其中 $\bm{B}_{,x}$ 表示对 $\bm{B}$ 再对 $x$ 求一次导数。

\textbf{4. 对两节点 Hermite 梁单元，把它写得更具体一点：}
我们前面已经得到
\[
\bm{B} =
\begin{bmatrix}
\dfrac{-6 + 12\xi}{L_e^2} &
\dfrac{-4 + 6\xi}{L_e} &
\dfrac{6 - 12\xi}{L_e^2} &
\dfrac{-2 + 6\xi}{L_e}
\end{bmatrix}.
\]
因此单元弯矩恢复公式为
\[
M^h(x)
=
EI
\begin{bmatrix}
\dfrac{-6 + 12\xi}{L_e^2} &
\dfrac{-4 + 6\xi}{L_e} &
\dfrac{6 - 12\xi}{L_e^2} &
\dfrac{-2 + 6\xi}{L_e}
\end{bmatrix}
\begin{bmatrix}
w_1 \\ \theta_1 \\ w_2 \\ \theta_2
\end{bmatrix}.
\]
由于这一行关于 $\xi$ 是\textbf{一次函数}，所以：
\begin{itemize}
    \item 两节点 Euler-Bernoulli 梁单元中的弯矩 $M^h(x)$ 在单元内部是\textbf{线性变化}的；
    \item 相应的剪力 $V^h(x)=M^{h\,\prime}(x)$ 在单元内部是\textbf{常数}。
\end{itemize}

\textbf{5. 为什么这件事值得特别警惕？}
如果真实梁承受均布载荷 $q_0$，解析解中：
\begin{itemize}
    \item 剪力一般是线性变化的；
    \item 弯矩一般是二次变化的。
\end{itemize}
但一个两节点 Hermite 梁单元恢复出来的结果却是：
\begin{itemize}
    \item 剪力常数；
    \item 弯矩线性。
\end{itemize}
这说明什么？说明\textbf{位移场即便足够准确，内力场仍然只是离散空间允许范围内的近似。}

\textbf{6. 直接微分恢复的优点与缺点：}
\begin{itemize}
    \item \textbf{优点：} 简单、直接、实现成本低。在任意采样点上都能算。
    \item \textbf{缺点：} 剪力比弯矩更敏感，因为它比位移多做了一次导数，数值误差会被进一步放大。
    \item \textbf{工程直觉：} 一般来说，位移最平滑、最稳定；转角次之；弯矩再次之；剪力最“躁动”，最容易出现跳变或噪声。
\end{itemize}

\textbf{7. 更稳妥的做法之一：恢复单元端内力向量}
除了在单元内部直接求导，你还可以从\textbf{单元内力平衡}的角度恢复端弯矩与端剪力。定义单元端内力向量
\[
\bm{f}_{\text{int}}^e = \bm{K}^e \bm{d}^e - \bm{F}_q^e.
\]
这个向量反映的是：为了平衡该单元的变形与分布载荷，单元两端需要提供怎样的剪力和弯矩。

通常可以把它理解为形如
\[
\bm{f}_{\text{int}}^e =
\begin{bmatrix}
V_1 \\ M_1 \\ V_2 \\ M_2
\end{bmatrix}
\]
的端截面内力结果（具体正负号要与你的自由度、节点力方向定义保持完全一致）。

\textbf{为什么这种恢复很重要？}
\begin{itemize}
    \item 它与单元刚度方程严格平衡；
    \item 用来画弯矩图、剪力图的端点值通常比“纯求导”结果更稳；
    \item 做结构后处理时，端内力往往是最先需要核对的量。
\end{itemize}

\textbf{8. 再进一步：用平衡关系恢复单元内部剪力与弯矩}
如果你已经通过端内力恢复得到了某一端截面的 $V$ 和 $M$，并且已知单元内分布载荷 $q(x)$，那么还可以利用静力平衡关系恢复内部内力分布。按本节约定，
\[
V'(x)=q(x), \qquad M'(x)=V(x).
\]
因此只要给定一个参考点（比如单元左端）的内力值，就可积分得到：
\[
V(x)=V(x_0)+\int_{x_0}^{x} q(s)\, ds,
\]
\[
M(x)=M(x_0)+\int_{x_0}^{x} V(s)\, ds.
\]
这种“平衡恢复”方法有一个很大的优点：\textbf{它强制满足单元内静力平衡}。在分布载荷已知时，它往往比简单对位移三次求导得到的剪力更符合工程直觉。

\textbf{9. 为什么剪力恢复通常比弯矩恢复更难？}
\begin{itemize}
    \item 弯矩只需要位移场的二阶导数；
    \item 剪力需要三阶导数，误差放大更明显；
    \item 多单元拼接后，弯矩通常可做到单元内平滑变化，而剪力常在单元边界出现分段常值或分段线性特征；
    \item 所以很多商业软件在显示剪力图时，本质上都做了某种平滑、投影或平衡修正。
\end{itemize}

\textbf{10. 对前面那个悬臂梁单元例子，后处理怎么做？}
前面我们已经求得
\[
\bm{d}^e =
\begin{bmatrix}
0 \\ 0 \\ \dfrac{q_0L^4}{8EI} \\ \dfrac{q_0L^3}{6EI}
\end{bmatrix}.
\]
将它代入弯矩恢复公式：
\[
M^h(x)=EI\,\bm{B}(x)\bm{d}^e.
\]
你会得到一个关于 $x$ 的\textbf{线性弯矩函数}；继续对它求导，就会得到一个\textbf{常值剪力}。这一点与前面从插值阶次得到的判断完全一致：两节点 Hermite 梁单元恢复出的弯矩在单元内是线性的，剪力在单元内是常数。

\textbf{这个结果要怎么理解？}
\begin{itemize}
    \item 它在\textbf{有限元离散意义下}是自洽的；
    \item 它能正确反映左端和右端截面内力的总体趋势；
    \item 但它并不是均布载荷悬臂梁解析解中的精确整段内力分布；
    \item 真正的解析解中，剪力沿梁长线性变化，弯矩沿梁长二次变化。
\end{itemize}
所以你应当把这个恢复结果理解为：\textbf{它是当前离散空间能表达出来的最佳单元内近似，而不是“真实连续解被完整还原”。}

\textbf{11. 工程上如何提高弯矩和剪力恢复质量？}
\begin{itemize}
    \item \textbf{细化网格：} 这是最稳定也最通用的方法；
    \item \textbf{使用高阶梁单元：} 提高位移插值阶次，内力恢复会更自然；
    \item \textbf{采用平衡恢复/投影恢复：} 让后处理结果同时兼顾位移兼容与静力平衡；
    \item \textbf{在积分点恢复再外推到节点：} 这是有限元后处理里非常常见的一条路；
    \item \textbf{避免直接把剪力图看得过“真”：} 尤其在粗网格下，剪力图更多是诊断趋势，而不是高精度真值。
\end{itemize}

\textbf{12. 这一小节真正想让你建立的意识是：}
\begin{itemize}
    \item 求解器解出来的是位移自由度，不是弯矩和剪力；
    \item 弯矩和剪力是\textbf{派生量}；
    \item 派生量的精度通常低于主变量位移；
    \item 后处理不是“画图步骤”，而是有限元结果可信度的重要组成部分。
\end{itemize}

\subsection{与前文的统一视角}
现在你应该已经能感受到，泊松方程和欧拉-伯努利梁其实在逻辑上高度统一：
\begin{itemize}
    \item 都从强形式方程出发；
    \item 都通过加权余量法与分部积分构造弱形式；
    \item 都会自然产生边界项，从而识别自然边界条件；
    \item 都通过 Galerkin 思想离散为矩阵方程；
    \item 只不过梁问题更进一步，暴露了\textbf{高阶问题对连续性的额外要求}。
\end{itemize}

\textbf{欧拉-伯努利梁这一节最值得你牢牢记住的三件事：}
\begin{enumerate}
    \item 四阶问题需要\textbf{两次}分部积分；
    \item 弱形式中出现 $w''$，所以标准位移型单元必须满足\textbf{$C^1$ 连续性}；
    \item Hermite 梁单元的节点自由度不是“只有位移”，而是\textbf{位移 + 转角}。
\end{enumerate}

\textbf{如果你能把这套推导完整复述出来，那么你对“为什么梁单元和实体单元不能随便混着理解”这件事，才算真正入门。}

\section{第九章：完整实战（三）—— Kirchhoff 薄板}
如果说欧拉-伯努利梁是一维四阶问题，那么 Kirchhoff 薄板就是它在二维空间中的“升级大 Boss”。梁只有一个空间坐标 $x$，而薄板有两个平面坐标 $(x,y)$；梁要求 $C^1$ 连续已经够麻烦了，薄板则要求\textbf{二维意义下的 $C^1$ 连续}，难度会陡然上升。

但也正因为如此，Kirchhoff 薄板是理解高阶连续介质有限元的绝佳样板。只要你把这一节吃透，就会真正明白：
\begin{itemize}
    \item 四阶问题为什么必须做两次分部积分；
    \item 为什么二维板单元远比一维梁单元难构造；
    \item 为什么很多工程程序宁愿转去用 Reissner-Mindlin 板，也不愿直接上 Kirchhoff 板。
\end{itemize}

下面我们从最经典的各向同性薄板弯曲问题出发，把强形式、弯矩-曲率关系、弱形式、离散化、边界条件、组装、后处理和工程实现一次走完整。

\textbf{符号约定：}
考虑位于中面区域 $\Omega \subset \mathbb{R}^2$ 上的一块薄板，横向挠度记为 $w(x,y)$，面外分布载荷记为 $p(x,y)$，板厚为 $h$，弹性模量为 $E$，泊松比为 $\nu$。对各向同性 Kirchhoff 薄板，其弯曲刚度定义为
\[
D = \frac{Eh^3}{12(1-\nu^2)}.
\]
为了让公式结构更清楚，本节把符号约定尽量写成矩阵形式。至于不同教材中弯矩正负号略有不同的问题，你只要在整套推导里始终保持同一约定，最终程序就不会乱。

\subsection{Kirchhoff 假设}
Kirchhoff 薄板理论的核心假设与欧拉-伯努利梁完全同源，只是从线扩展到了面：
\begin{itemize}
    \item 原来垂直于中面的法线，在变形后仍保持直线；
    \item 这条法线仍垂直于变形后的中面；
    \item 法线长度不变，因此横向剪切变形被忽略。
\end{itemize}

这三个假设的直接后果是：\textbf{薄板的弯曲能量只与挠度 $w$ 的二阶导数有关}，而不会像 Reissner-Mindlin 理论那样显式引入横向剪切应变。

\subsection{曲率、弯矩与剪力}
薄板有两个平面方向，因此“曲率”不再是梁里那一个 $w''$，而是一个曲率向量：
\[
\bm{\kappa} =
\begin{bmatrix}
w_{,xx} \\
w_{,yy} \\
2w_{,xy}
\end{bmatrix}.
\]
这里前两个分量是 $x$、$y$ 方向的弯曲曲率，第三个分量对应扭曲曲率。

对各向同性板，弯矩合力向量记作
\[
\bm{m} =
\begin{bmatrix}
M_x \\
M_y \\
M_{xy}
\end{bmatrix},
\]
其本构关系写为
\[
\bm{m} = \bm{D}_b \bm{\kappa},
\]
其中弯曲本构矩阵为
\[
\bm{D}_b
=
D
\begin{bmatrix}
1 & \nu & 0 \\
\nu & 1 & 0 \\
0 & 0 & \dfrac{1-\nu}{2}
\end{bmatrix}.
\]

\textbf{这组公式有非常重要的结构信息：}
\begin{itemize}
    \item $M_x$ 不只跟 $w_{,xx}$ 有关，也会耦合 $w_{,yy}$；
    \item $M_y$ 同理；
    \item $M_{xy}$ 由混合二阶导数 $w_{,xy}$ 控制；
    \item 这说明板弯曲是\textbf{二维耦合问题}，不像梁那样只有单一方向弯曲。
\end{itemize}

进一步定义广义剪力合力：
\[
Q_x = M_{x,x} + M_{xy,y},
\qquad
Q_y = M_{xy,x} + M_{y,y}.
\]
它们就是薄板内部由弯矩场“传递”出来的横向剪力结果量。

\subsection{强形式平衡方程}
板的局部平衡要求横向剪力继续满足散度平衡：
\[
Q_{x,x} + Q_{y,y} = p.
\]
把 $Q_x,Q_y$ 的定义代入，可得
\[
M_{x,xx} + 2M_{xy,xy} + M_{y,yy} = p.
\]
再将各向同性本构关系代入，可进一步化为著名的双调和方程
\[
D \nabla^4 w = p,
\]
其中
\[
\nabla^4 w = w_{,xxxx} + 2w_{,xxyy} + w_{,yyyy}.
\]

\textbf{这就是 Kirchhoff 薄板的强形式。}
请你特别注意：这里出现的是 $w$ 的\textbf{四阶偏导数}，所以这和梁一样，本质上是一个四阶问题；只不过梁是一维四阶，板是二维四阶。

\subsection{边界量与边界条件}
与梁问题类似，板边界上也会同时出现“位移型量”和“力学型量”，只是现在都沿着边界曲线定义。

设边界上一点的单位外法向为 $\bm{n}$，单位切向为 $\bm{s}$。则在边界局部坐标系下，可定义：
\begin{itemize}
    \item \textbf{法向转角：} $\dfrac{\partial w}{\partial n}$；
    \item \textbf{法向弯矩：} $M_{nn}$；
    \item \textbf{扭矩：} $M_{ns}$；
    \item \textbf{法向等效剪力：} $\widehat{Q}_n$。
\end{itemize}

其中
\[
M_{nn} = \bm{n}^T \bm{M}\bm{n},
\qquad
M_{ns} = \bm{n}^T \bm{M}\bm{s},
\]
这里 $\bm{M}$ 是由 $M_x,M_y,M_{xy}$ 组成的弯矩张量。常用的等效剪力定义为
\[
\widehat{Q}_n = Q_n + \frac{\partial M_{ns}}{\partial s}.
\]

\textbf{为什么要引入“等效剪力”这个看起来有点怪的量？}
因为对二维四阶问题做两次积分分部以后，边界上不仅会冒出剪力项，还会冒出扭矩沿边界切向变化的项。把它们整理到一起后，自然就形成了这个广义的法向等效剪力。

\textbf{Kirchhoff 板常见边界条件：}
\begin{itemize}
    \item \textbf{固支边 (Clamped Edge)：} $w=0,\; \dfrac{\partial w}{\partial n}=0$；
    \item \textbf{简支边 (Simply Supported Edge)：} $w=0,\; M_{nn}=0$；
    \item \textbf{自由边 (Free Edge)：} $M_{nn}=0,\; \widehat{Q}_n=0$。
\end{itemize}

\textbf{这一段一定要和梁类比着记：}
\begin{itemize}
    \item 梁里是“挠度 + 转角”对应本质边界条件；
    \item 板里则是“挠度 + 法向转角”对应本质边界条件；
    \item 梁里自然边界量是“弯矩 + 剪力”；
    \item 板里自然边界量是“法向弯矩 + 等效法向剪力”。
\end{itemize}

\subsection{弱形式构造}
如果你硬要从强形式 $D\nabla^4 w = p$ 出发逐项做两次二维分部积分，会得到一长串边界项，推起来非常累。但其最终结果，与从\textbf{总势能最小原理}或\textbf{虚功原理}得到的弱形式完全等价。

\textbf{1. 内部弯曲应变能：}
\[
U =
\frac{1}{2}
\int_{\Omega}
\bm{\kappa}^T \bm{D}_b \bm{\kappa}
\,
d\Omega.
\]

\textbf{2. 外载势能：}
\[
W =
\int_{\Omega} p\, w\, d\Omega
+
\int_{\Gamma_Q} \overline{\widehat{Q}}_n\, w\, d\Gamma
+
\int_{\Gamma_M} \overline{M}_{nn}\, \frac{\partial w}{\partial n}\, d\Gamma.
\]
如果边界上还存在角点集中力或角点集中力矩，那么还会出现附加的点载荷项。大多数基础教材先忽略角点项，本节也采用这个做法。

\textbf{3. 令总势能 $\Pi = U-W$ 一阶变分为零：}
\[
\delta \Pi = 0.
\]
于是得到 Kirchhoff 板的弱形式：
\[
\int_{\Omega}
\delta \bm{\kappa}^T \bm{D}_b \bm{\kappa}
\,
d\Omega
=
\int_{\Omega} p\, \delta w\, d\Omega
+
\int_{\Gamma_Q} \overline{\widehat{Q}}_n\, \delta w\, d\Gamma
+
\int_{\Gamma_M} \overline{M}_{nn}\, \delta \left(\frac{\partial w}{\partial n}\right) d\Gamma.
\]

\textbf{把它翻译成人话：}
\begin{itemize}
    \item 左边是板内部弯矩在虚曲率上所做的虚功；
    \item 右边是面外分布载荷、边界剪力和边界弯矩在虚位移与虚转角上所做的虚功；
    \item 这和梁的虚功原理是同一个家族，只是从一维推广到了二维。
\end{itemize}

\subsection{连续性要求与 $C^1$ 条件}
观察弱形式左端：
\[
\int_{\Omega}
\delta \bm{\kappa}^T \bm{D}_b \bm{\kappa}
\,
d\Omega.
\]
而曲率 $\bm{\kappa}$ 由 $w$ 的二阶导数组成，因此若使用标准位移型有限元，离散函数至少要能提供平方可积的二阶导数。这意味着：
\begin{itemize}
    \item $w$ 在单元交界处必须连续；
    \item $w_{,x}$ 与 $w_{,y}$ 也必须连续；
    \item 所以普通的 $C^0$ 拉格朗日四边形或三角形形函数\textbf{不够用}；
    \item 必须采用二维意义下的 $C^1$ 连续插值，或者使用非协调/混合/降阶理论替代。
\end{itemize}

\textbf{这就是 Kirchhoff 板比梁更难的根本原因。}
梁只要求一条线上的 $C^1$ 连续；板却要求一个二维网格上 $w,w_{,x},w_{,y}$ 在单元拼接处协调一致，构造难度要大得多。

\subsection{单元构造难点}
你现在已经可以理解，为什么工程里直接构造 Kirchhoff 板单元会比较“贵”：
\begin{itemize}
    \item \textbf{矩形相容元：} 例如 Bogner-Fox-Schmit (BFS) 这类双三次矩形元，通常每节点需要较丰富的自由度；
    \item \textbf{三角形相容元：} 例如 Argyris 元、HCT 元等，自由度设计更复杂；
    \item \textbf{非协调/退化方法：} 工程程序常转向 DKT、DKQ、MITC 或 Reissner-Mindlin 路线，再通过薄板极限逼近 Kirchhoff 行为。
\end{itemize}

\textbf{所以这一节最重要的不是背具体某一个板单元，而是抓住通用离散结构：}
\[
w^h = \bm{N}\bm{d}^e,
\qquad
\bm{\kappa}^h = \bm{B}\bm{d}^e.
\]
一旦这个结构你看懂了，具体板单元只是“形函数和自由度设计不同”。

\subsection{通用离散形式与板弯曲矩阵}
设某个 Kirchhoff 板单元的自由度向量为 $\bm{d}^e$，单元挠度插值写为
\[
w^h(x,y)=\bm{N}(x,y)\bm{d}^e.
\]
则曲率向量可写成
\[
\bm{\kappa}^h =
\begin{bmatrix}
w^h_{,xx} \\
w^h_{,yy} \\
2w^h_{,xy}
\end{bmatrix}
=
\bm{B}\bm{d}^e,
\]
其中
\[
\bm{B}
=
\begin{bmatrix}
N_{1,xx} & N_{2,xx} & \cdots \\
N_{1,yy} & N_{2,yy} & \cdots \\
2N_{1,xy} & 2N_{2,xy} & \cdots
\end{bmatrix}.
\]
注意这里的 $\bm{B}$ 与实体单元中的应变矩阵在结构上很像，但它包含的是\textbf{二阶导数}，这再次体现出 Kirchhoff 板对插值函数要求更高。

\subsection{单元刚度矩阵}
将离散场代入弱形式左端，可以得到 Kirchhoff 板单元的标准刚度阵公式：
\[
\bm{K}^e
=
\int_{\Omega^e}
\bm{B}^T \bm{D}_b \bm{B}
\,
d\Omega.
\]

\textbf{这条公式和前面的梁、泊松方程非常像，但层次更高：}
\begin{itemize}
    \item 泊松方程：$\bm{K}^e = \int \nabla N_i \cdot k \nabla N_j$；
    \item 欧拉-伯努利梁：$\bm{K}^e = \int EI\, \bm{B}^T\bm{B}$；
    \item Kirchhoff 板：$\bm{K}^e = \int \bm{B}^T \bm{D}_b \bm{B}$。
\end{itemize}
三者其实是同一套变分-离散逻辑在不同阶次、不同维度上的表现。

\subsection{单元载荷向量}
外载项同样可以离散为单元载荷向量。若单元上承受分布面载荷 $p(x,y)$，则一致载荷向量为
\[
\bm{F}^e_p
=
\int_{\Omega^e}
\bm{N}^T p
\,
d\Omega.
\]
若边界上存在等效剪力边载荷 $\overline{\widehat{Q}}_n$，则边界载荷项为
\[
\bm{F}^e_Q
=
\int_{\Gamma_Q^e}
\bm{N}^T \overline{\widehat{Q}}_n
\,
d\Gamma.
\]
若边界上存在法向弯矩边载荷 $\overline{M}_{nn}$，由于它作用在法向转角上，因此对应的离散形式不是简单的 $\bm{N}^T$，而是与 $\partial w/\partial n$ 的插值相关，通常可写成
\[
\bm{F}^e_M
=
\int_{\Gamma_M^e}
\bm{N}_{\theta n}^T \overline{M}_{nn}
\,
d\Gamma,
\]
其中 $\bm{N}_{\theta n}$ 表示法向转角的边界插值矩阵。

于是单元总载荷向量结构上就是
\[
\bm{F}^e = \bm{F}^e_p + \bm{F}^e_Q + \bm{F}^e_M.
\]

\textbf{这里一定要建立一个直觉：}
\begin{itemize}
    \item 面载荷通过挠度自由度做功；
    \item 边界弯矩通过法向转角自由度做功；
    \item 所以板单元的外载向量，往往比梁和泊松问题更“多层次”。
\end{itemize}

\subsection{参考单元、数值积分与实现}
如果你真的开始写代码，会发现 Kirchhoff 板单元的手推刚度阵往往长得吓人，因此实际实现通常依赖数值积分。

\textbf{1. 参考单元映射：}
和第五章介绍的通用离散化思想一样，板单元也会映射到标准参考单元上。比如矩形单元会映射到 $(\xi,\eta)\in[-1,1]^2$。

\textbf{2. 导数变换：}
由于曲率矩阵涉及二阶导数，除了普通的一阶 Jacobian 变换外，还要仔细处理二阶导数在物理坐标与参考坐标之间的转换。这比实体单元中一阶梯度变换更复杂，也更容易在程序里写错。

\textbf{3. 数值积分：}
Kirchhoff 板单元刚度阵依赖二阶导数乘积，积分多项式阶次较高，因此往往需要比普通低阶实体元更精细的积分方案。

\textbf{4. 一个非常实用的实现提醒：}
\begin{itemize}
    \item 先验证形函数是否满足节点插值与 $C^1$ 连续要求；
    \item 再验证 $\bm{B}$ 矩阵中的二阶导数是否正确；
    \item 再检查刚度阵是否对称、是否正定（在去掉刚体模态后）；
    \item 最后才去看位移和内力结果。
\end{itemize}
Kirchhoff 板元一旦写错，通常不是“结果偏一点”，而是刚度完全失真，甚至直接不收敛。

\subsection{全局组装与边界条件}
和所有有限元问题一样，单元层面的结果最终都要组装成全局系统：
\[
\bm{K}\bm{d} = \bm{F}.
\]
只不过这次全局自由度通常至少包含：
\begin{itemize}
    \item 节点挠度 $w$；
    \item 节点转角 $\theta_x = w_{,x}$；
    \item 节点转角 $\theta_y = w_{,y}$；
    \item 某些相容高阶矩形元还可能显式带有额外导数自由度。
\end{itemize}

\textbf{边界条件施加逻辑：}
\begin{itemize}
    \item 固支边：约束 $w$ 与 $\partial w/\partial n$；
    \item 简支边：约束 $w$，而 $M_{nn}=0$ 通过自然边界条件体现；
    \item 自由边：不施加位移型约束，自然边界量由右端项处理；
    \item 若边界方向与全局 $x,y$ 轴不重合，还需要把法向转角与全局转角自由度做坐标投影。
\end{itemize}

\textbf{这最后一句很关键：}
二维板边界条件比一维梁更容易出错，往往不是公式错，而是\textbf{局部法向/切向量和全局自由度方向没对应好}。

\subsection{曲率、弯矩与剪力后处理}
和梁问题完全一样，主求解器解出来的首先是位移型自由度，而不是内力结果。因此 Kirchhoff 板也必须做后处理恢复。

\textbf{1. 曲率恢复：}
\[
\bm{\kappa}^h = \bm{B}\bm{d}^e.
\]

\textbf{2. 弯矩恢复：}
\[
\bm{m}^h = \bm{D}_b \bm{\kappa}^h.
\]
也就是
\[
\begin{bmatrix}
M_x \\ M_y \\ M_{xy}
\end{bmatrix}
=
\bm{D}_b
\begin{bmatrix}
w_{,xx} \\ w_{,yy} \\ 2w_{,xy}
\end{bmatrix}.
\]

\textbf{3. 剪力恢复：}
\[
Q_x = M_{x,x} + M_{xy,y},
\qquad
Q_y = M_{xy,x} + M_{y,y}.
\]
你会立刻发现，剪力恢复比弯矩恢复更敏感，因为它需要对弯矩再做一次求导，也就是对位移场做\textbf{三阶导数}层级的恢复。和梁里“剪力比弯矩更躁动”的现象完全一样，只是在二维里更明显。

\textbf{4. 工程上通常怎么做：}
\begin{itemize}
    \item 在高斯积分点先恢复曲率与弯矩；
    \item 再通过外推、投影或平滑技术恢复到节点；
    \item 剪力常结合单元平衡关系单独恢复，而不是简单直接求导。
\end{itemize}

\subsection{标准算例工作流}
考虑一块长为 $a$、宽为 $b$ 的各向同性矩形薄板，承受均布面载荷 $p_0$。如果要用 Kirchhoff 板有限元完整求解，标准流程如下：
\begin{enumerate}
    \item 选定板单元类型，例如某个相容的 $C^1$ 矩形元或三角形元；
    \item 对区域 $\Omega=[0,a]\times[0,b]$ 划分网格；
    \item 在每个单元上建立挠度插值函数 $w^h=\bm{N}\bm{d}^e$；
    \item 计算曲率矩阵 $\bm{B}$ 与单元刚度阵 $\bm{K}^e=\int \bm{B}^T\bm{D}_b\bm{B}$；
    \item 计算均布面载荷对应的一致载荷向量 $\bm{F}^e_p=\int \bm{N}^Tp_0$；
    \item 将所有单元刚度阵与载荷向量组装到全局系统；
    \item 根据边界类型施加固支、简支或自由边条件；
    \item 解线性方程组得到全局自由度 $\bm{d}$；
    \item 在单元积分点恢复 $\bm{\kappa},\bm{m},\bm{Q}$；
    \item 绘制挠度云图、弯矩等值线以及边界反力。
\end{enumerate}

\textbf{你会发现，流程本身和梁、泊松方程一模一样。}
真正增加难度的不是“有限元框架变了”，而是：
\begin{itemize}
    \item 方程阶次更高；
    \item 变量维度从一维升到二维；
    \item 连续性要求从 $C^0$ 升到二维 $C^1$；
    \item 单元设计与边界条件实现都更复杂。
\end{itemize}

\subsection{与前文的统一视角}
现在可以把这一整章的三个代表性问题并排看：
\begin{itemize}
    \item \textbf{泊松方程：} 二阶标量问题，弱形式后只需一阶导，标准 $C^0$ 单元即可；
    \item \textbf{欧拉-伯努利梁：} 一维四阶问题，弱形式后仍需二阶导，因此要求一维 $C^1$；
    \item \textbf{Kirchhoff 薄板：} 二维四阶问题，弱形式后仍需二阶导，因此要求二维 $C^1$。
\end{itemize}

\textbf{这三者共同构成了一条非常清晰的有限元认知链：}
\begin{enumerate}
    \item 看强形式的最高阶导数；
    \item 看分部积分后弱形式里还剩几阶导数；
    \item 由此判断离散空间至少需要什么连续性；
    \item 再决定该选哪种形函数、哪种单元、哪种实现路线。
\end{enumerate}

\textbf{Kirchhoff 板这一节最重要的结论有四个：}
\begin{enumerate}
    \item 薄板弯曲本质上是二维四阶问题；
    \item 弯曲能量由二阶导数组成的曲率控制；
    \item 标准位移型 Kirchhoff 板有限元要求二维 $C^1$ 连续；
    \item 这正是它在理论上优雅、在实现上困难的根源。
\end{enumerate}

\textbf{如果你能把梁和 Kirchhoff 板放在同一个逻辑框架下去理解，那么你对高阶有限元的骨架，基本就真正搭起来了。}

\section{第十章：结语}
前六章给出了有限元理论的骨架：从数学符号、散度定理、分部积分，到加权余量法、变分法、函数空间、离散化理论与稳定性分析；第七至第九章则把这些理论分别投射到二阶标量问题、一维四阶问题和二维四阶问题上，形成了完整的推导闭环。

\textit{“代码是理论的影子，理论是代码的灵魂。当你推导出的每一个积分项都能在你的 Subroutine 中找到对应的循环时，你便真正掌握了有限元。”}

\end{document}
